\chapter{Token Type Registry and Standard Token Types}\label{app:token-registry}

This appendix specifies the registries that bind token type identifiers (Sec.~\ref{sec:token-types}) to their definitions, and catalogues the standard mint-reason formats, the issuance policies composed from them, and the standard token types. The appendix is developed iteratively: items advance from outline through draft to specified across document revisions, with review feedback in between. The specification roadmap and status is in Sec.~\ref{app:staging}.

\section{Registries}\label{app:registries}

\index{registry!token type}\index{Unicity Foundation}
The Unicity Foundation maintains machine-readable identifier registries in a public repository\footnote{\url{https://github.com/unicitynetwork/unicity-ids}}. Three registries are relevant to tokens:

\begin{description}
\item[Semantic tags] (\texttt{cbor-tags.json}) Global CBOR tag assignments, Appendix~\ref{app:cbor-tags}. Mint-reason formats are keyed by these tags.
\item[Asset metadata] (\texttt{unicity-ids.\textit{network}.json}) Display metadata of deployed assets, keyed by token type identifier; schema in Sec.~\ref{app:asset-metadata}. This is the currently operational registry; it covers external metadata only.
\item[Token type definitions] The registry of (type identifier, definition manifest) pairs specified in Sec.~\ref{app:type-registry}. It extends the asset-metadata registry with the consensus-relevant fields of $\mathcal{D}_\mathsf{ty}$ (Sec.~\ref{sec:token-type}): admissible reason formats, payload rules and issuance policy parameters.
\end{description}

Registry entries are coordination data, not trust roots. The consensus-relevant binding between a type identifier and its rules comes from the frozen definition itself --- self-authenticating when the identifier is content-addressed (Sec.~\ref{sec:token-type}) --- while the registry provides discovery, naming, and the social agreement on which definitions exist. Nevertheless, consensus-relevant registry fields MUST be append-only: an entry, once published as active, is never modified or reassigned, only marked deprecated.

\section{Asset Metadata Schema}\label{app:asset-metadata}

\index{registry!asset metadata}
An asset metadata entry is a JSON object with the following fields:

\begin{longtable}{llp{7.6cm}}
\textbf{Field} & \textbf{Type} & \textbf{Description} \\
\hline
\endhead
\texttt{network} & string & network name, e.g.\ \texttt{"unicity:testnet2"} \\
\texttt{assetKind} & string & \texttt{"fungible"} or \texttt{"non-fungible"} \\
\texttt{name} & string & human-readable asset name \\
\texttt{symbol} & string & ticker symbol (fungible assets) \\
\texttt{decimals} & integer & display scaling factor (fungible assets) \\
\texttt{description} & string & free-text description \\
\texttt{icons} & array & objects with an \texttt{url} field, pointing to icon images \\
\texttt{id} & string & lowercase hexadecimal token type identifier $\mathsf{ty}$ \\
\end{longtable}

Example (Unicity testnet v2 native coin):
\begin{quote}\small\ttfamily
\{ "network": "unicity:testnet2", "assetKind": "fungible",\\
\phantom{\{ }"name": "unicity", "symbol": "UCT", "decimals": 18,\\
\phantom{\{ }"description": "Unicity testnet v2 native coin",\\
\phantom{\{ }"icons": [\{"url": "https://\ldots/unicity\_logo\_32.png"\}],\\
\phantom{\{ }"id": "f581d30f593e4b369d684a4563b5246f07b1d265f7178a2c0a82b81f39c24dc0" \}
\end{quote}

All fields are presentation-level. In particular, \texttt{decimals} affects display only: on-ledger amounts $v$ are non-negative integers in minimal units (Sec.~\ref{sec:asset-structure}), and no consensus rule depends on the scaling factor. Wallets SHOULD refuse to render assets whose registry entry is missing, rather than display raw identifiers as if they were named assets --- the display-layer analogue of the fail-shut policy (Sec.~\ref{sec:mint-authorization}).

\section{Token Type Registry}\label{app:type-registry}

A token type registry entry publishes the definition manifest behind a type identifier. The manifest is the canonical CBOR structure $b_\mathcal{D}$ from which a content-addressed identifier is derived (Sec.~\ref{sec:token-type}); the registry stores it together with presentation metadata:

\begin{longtable}{lp{9.6cm}}
\textbf{Field} & \textbf{Description} \\
\hline
\endhead
$\mathsf{ty}$ & token type identifier ((forever) unique registry key) \\
status & \texttt{draft} $\mid$ \texttt{active} $\mid$ \texttt{deprecated} \\
standard & the standard token type this entry instantiates (e.g.\ UTS-1, Sec.~\ref{app:standard-types}), or a reference (URL and hash) to a frozen custom specification \\
issuance policy & the policy of Sec.~\ref{app:issuance-policies} and, for capped issuance, its profile \\
$\mathsf{RT}$ & admissible mint-reason format tags (Sec.~\ref{app:reason-formats}) \\
payload rules & parameters of $\Call{ParseAux}{}$, e.g.\ the permitted asset identifiers of a fungible type \\
issuance parameters & parameters of $\Call{VerifyIssuance}{}$, e.g.\ the issuer predicate hash $H(\mathsf{CBOR}(\predi_\mathsf{issue}))$, a bridge configuration hash $h_\mathsf{cfg}$ (Appendix~\ref{app:bridging-config}), or a proof-system verification key hash $h_\mathsf{vk}$ (Sec.~\ref{app:reason-proof}) \\
metadata & the asset metadata entry of Sec.~\ref{app:asset-metadata} \\
\end{longtable}

Lifecycle rules:
\begin{itemize}
	\item An entry becomes \texttt{active} only with a complete, frozen definition: allocated reason tags, published parsing and verification rules, and --- for proof-based types --- the frozen verification key.
	\item Once \texttt{active}, all fields except status and presentation metadata are immutable. Fixing or extending rules means registering a \emph{new} type (cf.\ the vault upgrade discipline, Sec.~\ref{sec:bridge-security}).
	\item Deprecation is advisory: tokens of a deprecated type remain verifiable; whether to keep the type in the local $\mathcal{TD}$ is each application's decision.
	\item For content-addressed identifiers, anyone can audit the registry by recomputing $\mathsf{ty} = H(\mathsf{TYPE\_DEF}, b_\mathcal{D})$; a mismatch proves the entry corrupt.
\end{itemize}

\section{Standard Mint-Reason Formats}\label{app:reason-formats}

\index{mint reason!formats}
The following reason formats are defined or planned; statuses are those of the specification text (Sec.~\ref{app:staging}).

\begin{longtable}{llll}
\textbf{Tag} & \textbf{Format} & \textbf{Status} & \textbf{Specification} \\
\hline
\endhead
$\bot$ & absent reason & specified & admissible iff $\bot \in \mathsf{RT}$ (Sec.~\ref{sec:mint-authorization}) \\
39042 & MintReasonGenesis & reserved & native-coin genesis allocation \\
39043 & MintReasonReward & reserved & Validator reward issuance \\
39044 & MintReasonSplit & specified & Sec.~\ref{sec:token-splitting}, Appendix~\ref{app:split-mint-reason} \\
39048 & BridgeMintReason & specified & Sec.~\ref{sec:bridge-in}, Appendix~\ref{app:bridging-backing-reason} \\
$\mathsf{ISSUER\_TAG}$ & MintReasonIssuer & draft & Sec.~\ref{app:reason-issuer} \\
$\mathsf{SLOT\_TAG}$ & MintReasonSlot & draft & Sec.~\ref{app:reason-slot} \\
$\mathsf{VRF\_TAG}$ & MintReasonVRF & draft & Sec.~\ref{app:reason-vrf} \\
$\mathsf{SERIAL\_TAG}$ & MintReasonSerial & draft & Sec.~\ref{app:reason-serial} \\
$\mathsf{SPLITREF\_TAG}$ & MintReasonSplitRef & draft & Sec.~\ref{app:reason-splitref} \\
$\mathsf{MERGE\_TAG}$ & MintReasonMerge & draft & Sec.~\ref{app:reason-merge} \\
$\mathsf{CONVERT\_TAG}$ & MintReasonConvert & draft & Sec.~\ref{app:reason-convert} \\
$\mathsf{PROOF\_TAG}$ & MintReasonProof & draft & Sec.~\ref{app:reason-proof} \\
$\mathsf{NAME\_TAG}$ & MintReasonName & idea & n/a \\
\end{longtable}

\noindent Tag values for the draft formats (including the auxiliary manifest tags defined below) are to be allocated in the semantic tag registry (Appendix~\ref{app:cbor-tags}) when the drafts are frozen.

\subsection{Genesis Commitment}\label{app:genesis-commitment}

The formats below authorize one specific mint by signing or proving a commitment over all mint fields \emph{except} the reason itself (excluded to avoid a circular commitment, exactly as in the split output commitment, Appendix~\ref{app:split-output-commitment}). Define $\mathsf{MINT\_AUTH}$ as the ASCII byte string \texttt{"UNICITY\_MINT\_AUTH"}. For mint data $D_\mathsf{mint} = (\alpha, \predi', \mathsf{salt}, \mathsf{ty}, \mathsf{reason}, \auxd')$, the \emph{genesis commitment} is
\[
m_\mathsf{mint} := \mathsf{SHA\text{-}256}\bigl(\mathsf{CBOR}(
	\mathsf{MINT\_AUTH},\;
	\alpha,\;
	\encpred(\predi'),\;
	\mathsf{salt},\;
	\mathsf{id},\;
	\mathsf{ty},\;
	\auxd'
)\bigr),
\]
with $\mathsf{id} = H(\mathsf{salt}, \alpha)$ and each term encoded as in Appendix~\ref{app:split-output-commitment}. The commitment binds an authorization to exactly one network, recipient, token identifier, token type and payload; an authorization produced for one mint cannot be replayed for any other. Re-presenting the same authorization can only re-create the same single token, whose genesis state is certified at most once (Sec.~\ref{sec:request-validation}).

\subsection{Issuer-Signature Reason}\label{app:reason-issuer}

\index{mint reason!issuer signature}
Minting is gated by an issuer's authorization. Rather than fixing a signature scheme, the format reuses the predicate framework (Sec.~\ref{sec:predicates}): the type definition pins an \emph{issuance predicate} $\predi_\mathsf{issue}$, and the reason carries an unlocking argument for it. Single-key signature, multi-signature, $k$-of-$n$ threshold, P2SH indirection and time-boxed issuance windows are thereby available with no new machinery.

\begin{itemize}
	\item \textbf{Type parameters}: $\predi_\mathsf{issue}$, stored in the definition manifest as the predicate hash $H(\mathsf{CBOR}(\predi_\mathsf{issue}))$ or the full encoded predicate.
	\item \textbf{Reason body}: CBOR array under tag $\mathsf{ISSUER\_TAG}$:
\end{itemize}

\begin{longtable}{clp{8.6cm}}
\textbf{Index} & \textbf{Type} & \textbf{Description} \\
\hline
0 & unsigned integer & format version; this specification defines version $1$ \\
1 & CBOR structure & unlocking argument $u$ for $\predi_\mathsf{issue}$, encoded as in Appendix~\ref{app:predicate-encoding} \\
\end{longtable}

\textbf{Verification} (as $V_{\mathsf{ISSUER\_TAG}}$ in Alg.~\ref{alg:mint-auth}):
\begin{enumerate}
	\item Decode $(\mathsf{ver}, u)$ from the reason body; ensure $\mathsf{ver} = 1$.
	\item $\predi_\mathsf{issue} \gets$ issuance predicate pinned by $\mathcal{D}_\mathsf{ty}$; if the manifest pins only the hash, the full predicate is carried in $u$ and checked against the hash (as in P2SH, Sec.~\ref{sec:p2sh}).
	\item $m_\mathsf{mint} \gets$ genesis commitment of $D_\mathsf{mint}$ (Sec.~\ref{app:genesis-commitment}).
	\item Ensure $\predi_\mathsf{issue}(\mathsf{time}(\pi_0), m_\mathsf{mint}, u) = 1$, with $\mathsf{time}(\pi_0)$ extracted from the genesis inclusion proof (Sec.~\ref{sec:time-extraction}).
\end{enumerate}

\paragraph{Security notes.}
\begin{description}
	\item[No replay] $m_\mathsf{mint}$ binds the authorization to one token identifier, recipient and payload (Sec.~\ref{app:genesis-commitment}).
	\item[No retroactive revocation] A distributed, offline-verifiable token cannot be invalidated after issuance; compromised issuer keys are handled by deprecating the type and registering a successor, which stops \emph{future} issuance only.
	\item[Supply discipline] This format alone leaves the number of authorized geneses unbounded. A type can nevertheless enforce a hard cap by additionally fixing its only permitted direct-issuance identity and amount (the singleton profile, Sec.~\ref{app:uts-singleton}). Independently issued capped series use the VRF or committed-slot formats below. The complete type policy, including every admitted reason format, determines whether issuance is bounded.
\end{description}

\subsection{Committed-Slot Reason}\label{app:reason-slot}

\index{mint reason!slot-bound}\index{issuance limit}
This optional format supports independent issuance with arbitrary pre-committed denominations or NFT content (profile S, Sec.~\ref{app:uts-slots}). The type definition commits a finite set of $N$ \emph{issuance slots}; each slot fixes one genesis certification state. The Unicity Service's single-spend rule (Sec.~\ref{sec:request-validation}) then permits at most one certified genesis per slot on a given network. A recipient checks its own commitment proof locally; no query for other tokens is needed. The singleton and VRF profiles avoid constructing this table when arbitrary per-slot commitments are unnecessary (Sec.~\ref{app:uts-capped}).

\subsubsection{Slot Commitment}

At type creation the issuer samples independent slot salts $\mathsf{salt}_1, \ldots, \mathsf{salt}_N$, each with at least 128 bits of min-entropy (a stronger requirement than the generic salt rule of Sec.~\ref{sec:mint-tx}: here unpredictability is security-relevant, see below), and fixes per-slot payload constraints $b_1, \ldots, b_N$ with $b_j \in \bytestr \opt$: a non-$\bot$ value mandates the exact genesis payload byte string of slot $j$ (e.g.\ its denomination); $\bot$ leaves the payload to mint time. Define the ASCII domain separators $\mathsf{SLOT\_SALT} = \texttt{"UNICITY\_SLOT\_SALT"}$ and $\mathsf{SLOT\_LEAF} = \texttt{"UNICITY\_SLOT\_LEAF"}$, and derive
\[
\begin{aligned}
h_j &:= \mathsf{SHA\text{-}256}(\mathsf{CBOR}(\mathsf{SLOT\_SALT}, \mathsf{salt}_j)),\\
\mathsf{leaf}_j &:= \mathsf{SHA\text{-}256}(\mathsf{CBOR}(\mathsf{SLOT\_LEAF}, h_j, b_j)).
\end{aligned}
\]
The type manifest pins the slot mode, $N$, $r_\mathsf{slots}$ and $\predi_\mathsf{issue}$, with $1 \le N \le 2^{32}$. Let $k_j$ be the 256-bit big-endian encoding of $j-1$. The mode fixes the tree and its interpretation; a mint cannot choose another mode:
\begin{description}
	\item[Count mode (NFTs)] $r_\mathsf{slots}$ is the root of an RSMT (Appendix~\ref{app:rsmt}) mapping $k_j$ to $\mathsf{leaf}_j$. The payload constraint $b_j$ is either an exact NFT payload or $\bot$.
	\item[Value mode (fungible assets)] The manifest additionally pins one asset identifier $\mathsf{aid}$ and $1 \le X < 2^{256}$. Every $b_j$ MUST be the canonical encoding of $\{(\mathsf{aid},v_j)\}$, with $v_j \ge 1$. Build an RSMST (Appendix~\ref{app:rsmst}) with leaves $(k_j,\mathsf{leaf}_j,v_j)$; its root hash is $r_\mathsf{slots}$ and its root sum MUST equal $X$. All amounts, subtree sums and the total are strictly below $2^{256}$; overflow rejects.
\end{description}

The issuer SHOULD distribute the \emph{audit table} $\langle(h_1, b_1), \ldots, (h_N, b_N)\rangle$ alongside the type definition when users need to inspect the entire allocation or content set. This is an immutable file, distributable by ordinary file hosting or with the definition; it needs no public blockchain or mutable shared registry. The table reveals no salt and lets readers reconstruct the committed tree. It is \emph{not required to verify the cap}: count mode bounds the accepted indices, and value mode proves the committed total from each token's sum proof. A plain Merkle membership proof would not establish that total. The salts themselves MUST remain confidential until their mints have been certified (Sec.~\ref{app:capped-availability}).

\subsubsection{Slot Uniqueness}

The mint of slot $j$ MUST use $D_\mathsf{mint}.\mathsf{salt} = \mathsf{salt}_j$. The genesis certification state is then a deterministic function of the slot: $\mathsf{id} = H(\mathsf{salt}_j, \alpha)$, source state hash $\sthash = H(\mathsf{id}, \mathsf{MINT\_SUFFIX})$, minting key $\mathsf{pk}_\mathsf{mint}$ derived from $\mathsf{id}$ (Appendix~\ref{app:el-constants}), and state identifier $\mathsf{sid} = H(\mathsf{pk}_\mathsf{mint}, \sthash)$. The Unicity Service certifies at most one transaction per state identifier (Sec.~\ref{sec:request-validation}); hence at most one certified genesis per slot can ever exist on network $\alpha$ --- a second mint of slot $j$, by the issuer or anyone else, is rejected exactly as a double spend would be.

Consequently, at most $N$ geneses can be created by slot issuance. In value mode the RSMST proof additionally bounds the sum of all provable denominations by $X$, including against a maliciously constructed tree (Sec.~\ref{sec:rsmst-properties}). Admitting only the slot format and the value-conserving split/merge formats bounds outstanding fungible value by $X$, but does not bound the number of split-created containers. In count mode, split and merge MUST be excluded, so at most $N$ NFT items can ever be issued. No additional direct-issuance format is admitted by this profile.

\subsubsection{Reason Body}

CBOR array under tag $\mathsf{SLOT\_TAG}$:

\begin{longtable}{clp{8.6cm}}
\textbf{Index} & \textbf{Type} & \textbf{Description} \\
\hline
0 & unsigned integer & format version; this specification defines version $1$ \\
1 & unsigned integer & slot index $j$, $1 \le j \le N$ \\
2 & CBOR array & count mode: RSMT inclusion certificate; value mode: RSMST proof (Sec.~\ref{app:split-allocation-proof}) \\
3 & byte string or \texttt{null} & the committed payload constraint $b_j$ \\
4 & Unlocking Argument & unlocking argument $u$ for $\predi_\mathsf{issue}$ (Appendix~\ref{app:predicate-encoding}); CBOR structure \\
\end{longtable}

\subsubsection{Verification}

As $V_{\mathsf{SLOT\_TAG}}$ in Alg.~\ref{alg:mint-auth}:
\begin{enumerate}
	\item Decode $(\mathsf{ver}, j, \pi_\mathsf{slot}, b_j, u)$; ensure $\mathsf{ver} = 1$ and $1 \le j \le N$.
	\item With $s=D_\mathsf{mint}.\mathsf{salt}$, recompute
\[
\begin{aligned}
h &\gets \mathsf{SHA\text{-}256}(\mathsf{CBOR}(\mathsf{SLOT\_SALT},s)),\\
\mathsf{leaf} &\gets \mathsf{SHA\text{-}256}(\mathsf{CBOR}(\mathsf{SLOT\_LEAF},h,b_j)).
\end{aligned}
\]
Let $k_j$ be the 256-bit big-endian encoding of $j-1$.
	\item In count mode, ensure $\textsc{rsmt\_verify\_inclusion}(\pi_\mathsf{slot}, r_\mathsf{slots}, k_j, \mathsf{leaf}) = 1$ (Sec.~\ref{sec:rsmt-verify-inclusion}). In value mode, require $b_j \ne \bot$ and decode it as exactly $\{(\mathsf{aid},v_j)\}$ for the pinned asset, with $v_j \ge 1$; verify the RSMST proof for $(k_j,\mathsf{leaf},v_j)$ against $r_\mathsf{slots}$ and require success and reconstructed root sum $X$ (Sec.~\ref{sec:rsmst-verify-inclusion}).
	\item If $b_j \ne \bot$: ensure $D_\mathsf{mint}.\auxd' = b_j$ byte-exactly.
	\item $m_\mathsf{mint} \gets$ genesis commitment (Sec.~\ref{app:genesis-commitment}); ensure $\predi_\mathsf{issue}(\mathsf{time}(\pi_0), m_\mathsf{mint}, u) = 1$.
\end{enumerate}
The driver additionally validates the payload against $\Call{ParseAux}{}$ and the issuance policy (Alg.~\ref{alg:mint-auth}), as for every format.

\paragraph{Security notes.}
\begin{description}
	\item[Hard limit] At most $N$ direct slot geneses per network; in value mode their cumulative value is at most the manifest's $X$, even if the issuer is malicious or its keys are compromised. Re-minting split/merge outputs conserves that value and is not additional direct issuance.
	\item[Protection against advance blocking] The audit table exposes $h_j$, not the salt or token identifier. Under the hash assumptions, the public commitment and certification request do not reveal the minting private key. Keeping the token data confidential until certification protects unused slots against advance blocking; the confidentiality and disclosure requirements are in Sec.~\ref{app:capped-availability}.
	\item[Separate custody] A valid mint requires both $\mathsf{salt}_j$ and an unlocking of $\predi_\mathsf{issue}$. With independent custody, compromise of either alone cannot authorize issuance. Leaked salts nevertheless permit permanent slot blocking; issuance safety and availability are different properties (Sec.~\ref{app:capped-availability}).
	\item[Fail-safe availability] Lost salts strand the remaining slots: value fails to be created but is never inflated or redirected. The issuer is responsible for durable, confidential storage of unused salts.
	\item[Denomination rigidity] Committed denominations are fixed at type creation; economic flexibility comes after minting, via split and merge, which conserve the capped total.
\end{description}

\paragraph{Open (claim-drop) variant.} As an explicit exception to the salt-confidentiality requirements, a deployment MAY publish the salts and set $\predi_\mathsf{issue}$ to the trivially-true predicate to permit public claims. The first \emph{valid} mint to consume a slot creates its token, but an arbitrary certification request can consume the slot without creating any valid token. This variant preserves the cap and deliberately accepts permanent blocking; it does not guarantee that every slot yields a successful claim.

\subsection{VRF Serial Reason}\label{app:reason-vrf}

\index{mint reason!VRF}\index{verifiable random function}
This format derives a unique mint salt for each serial in a frozen range, without generating or distributing a slot table. It is the default for capped NFT collections (profile V, Sec.~\ref{app:uts-vrf}) and also supports independent, equal-denomination fungible issuance. A verifiable random function (VRF) makes the derivation publicly checkable at mint time while keeping unused mint identities unpredictable to parties without the VRF secret key.

\subsubsection{Parameters and Derivation}

The manifest pins $(\mathsf{lo},\mathsf{hi},\mathsf{suite},\mathsf{pk}_\mathsf{vrf},\predi_\mathsf{issue})$, with $0 \le \mathsf{lo} \le \mathsf{hi} < 2^{32}$ and $N=\mathsf{hi}-\mathsf{lo}+1$. The RECOMMENDED instantiation is \texttt{ECVRF-EDWARDS25519-SHA512-ELL2} from RFC~9381, Sec.~5.5\footnote{\url{https://www.rfc-editor.org/rfc/rfc9381.html}}; this draft's version~1 fixes that suite (manifest value $\mathsf{suite}=4$, RFC suite byte \texttt{0x04}). It uses a 32-byte public key, an 80-byte proof and a 64-byte output. Implementations MUST use the RFC algorithms and encodings, including its key validation with \texttt{validate\_key = TRUE}, and SHOULD check interoperability against Appendix~B.4 of that RFC. Other suites require a separately specified format version.

Define the ASCII byte strings $\mathsf{VRF\_INPUT}=\texttt{"UNICITY\_VRF\_INPUT"}$ and $\mathsf{VRF\_SALT}=\texttt{"UNICITY\_VRF\_SALT"}$. For a serial $j$ in the range, compute
\[
\begin{aligned}
x_j &:= \mathsf{CBOR}(\mathsf{VRF\_INPUT},\alpha,\mathsf{ty},j),\\
\pi_j &:= \mathsf{ECVRF\_prove}(\mathsf{sk}_\mathsf{vrf},x_j),\\
y_j &:= \mathsf{ECVRF\_proof\_to\_hash}(\pi_j),\\
\mathsf{salt}_j &:= \mathsf{SHA\text{-}256}(\mathsf{CBOR}(\mathsf{VRF\_SALT},x_j,y_j)).
\end{aligned}
\]
Here each $\mathsf{CBOR}$ expression encodes one deterministic array; $\alpha$ and $j$ are unsigned integers, and the other terms are byte strings. The exact bytes $x_j$ are the VRF input, without an additional prehash. Salt derivation uses the \emph{verified VRF output}, never the proof bytes: different proofs for the same output MUST NOT create different mint identities. The manifest is frozen and $\mathsf{ty}$ is determined before these values are computed, so the derivation has no circular commitment. The issuer MUST NOT substitute another key, suite, range or input encoding for the same type.

\subsubsection{Reason Body and Verification}

The body under $\mathsf{VRF\_TAG}$ is the four-element CBOR array $[1,j,\pi_j,u]$: version and serial are unsigned integers, $\pi_j$ is an 80-byte string, and $u$ is an unlocking argument for the pinned issuance predicate. As $V_{\mathsf{VRF\_TAG}}$ in Alg.~\ref{alg:mint-auth}:
\begin{enumerate}
	\item Require the exact array shape, version $1$, $\mathsf{suite}=4$, the prescribed key/proof lengths, and $\mathsf{lo} \le j \le \mathsf{hi}$; reject malformed encodings.
	\item Recompute $x_j$ from the certified mint's network and type and the reason's serial. Run $\mathsf{ECVRF\_verify}(\mathsf{pk}_\mathsf{vrf},x_j,\pi_j)$ with \texttt{validate\_key = TRUE}. Reject every RFC decoding, key-validation or proof-validation failure; take $y_j$ only from a successful verification result.
	\item Recompute $\mathsf{salt}_j$ from $(x_j,y_j)$ and require $D_\mathsf{mint}.\mathsf{salt}=\mathsf{salt}_j$ byte-for-byte. The ordinary mint checks derive the token identifier and genesis certification state from that salt.
	\item Require $\predi_\mathsf{issue}(\mathsf{time}(\pi_0),m_\mathsf{mint},u)=1$, using the genesis commitment of Sec.~\ref{app:genesis-commitment}. The driver also applies the profile's payload and issuance checks, including the fixed denomination for fungible deployments.
\end{enumerate}

\paragraph{Security and operation.} The required property for the cap is one verifiable output per fixed key and input, even when the issuer chooses the key maliciously. The recommended ECVRF provides this full uniqueness property; key validation additionally provides collision resistance under malicious key generation (RFC~9381, Sec.~7.1.1). A generic signature hashed into a salt is not a substitute: a signer may be able to produce multiple valid signatures for one serial.

The cap survives disclosure of the VRF secret key, but unused states then become predictable and blockable. The issuer SHOULD generate and store the VRF key independently of its issuance-signing keys, and MUST keep unused salts, outputs and proofs confidential until certification. An exposed VRF proof reveals the output even without a valid issuer authorization. A public VRF-evaluation endpoint for unused serials therefore defeats the availability protection. Losing the VRF secret strands unused serials unless their proofs and salts were already backed up; changing keys requires a new type. Pseudorandomness assumes proper secret-key generation and secrecy and is not guaranteed against an issuer deliberately revealing its own key (RFC~9381, Sec.~7.1.2). These operational requirements protect availability, not the hard cap.

\subsection{Public Serial Reason}\label{app:reason-serial}

This optional format is for NFT collections that deliberately accept predictable mint states (profile P, Sec.~\ref{app:uts-public-serial}). The manifest pins $(\mathsf{lo},\mathsf{hi},\predi_\mathsf{issue})$ with the same range bounds as the VRF format. Define $\mathsf{SERIAL\_SALT}$ as the ASCII byte string \texttt{"UNICITY\_SERIAL\_SALT"}, and require
\[
\mathsf{salt}_j := \mathsf{SHA\text{-}256}\bigl(\mathsf{CBOR}(\mathsf{SERIAL\_SALT},\alpha,\mathsf{ty},j)\bigr).
\]
The expression encodes one deterministic array, with unsigned integer network and serial and byte-string domain separator and type. The reason body under $\mathsf{SERIAL\_TAG}$ is $[1,j,u]$. Verification requires this exact shape and version, an unsigned $j$ in the pinned range, byte-exact equality of the derived salt to the mint salt, and $\predi_\mathsf{issue}(\mathsf{time}(\pi_0),m_\mathsf{mint},u)=1$. Ordinary mint and NFT payload checks also apply. The serial is read from this verified reason; a display field cannot override it.

The bounded input range and fixed derivation permit at most $N$ NFT geneses per network. This is a cryptographic uniqueness claim under the hash assumptions, not an assertion that a finite hash is mathematically injective. The derivation includes the type and network to separate collections and networks, and excludes recipient, content, signatures and user-selected randomness so they cannot create a second identity for one serial. Anyone can nevertheless pre-consume these public states with arbitrary certification requests; see Sec.~\ref{app:capped-availability}. Issuer signatures checked by recipients do not prevent that attack.

\subsection{Reference-Form Split Reason}\label{app:reason-splitref}

\index{mint reason!split, reference form}\index{reference form}
The split reason of Sec.~\ref{sec:token-splitting} (tag 39044) embeds the complete certified burned source token in the mint reason, hence in $\txhash$, hence in the certified genesis. Those bytes are fixed forever: a token minted under that format can never have its provenance compressed (Sec.~\ref{sec:token-evidence}). This format is its reference-form counterpart --- identical in what it proves and in every conservation check, differing only in that the source is named by identifier and supplied as evidence beside the token.

It is a separate registry entry with its own tag, not a revision of 39044: the split reason is frozen as \emph{specified}, and the registry discipline (Sec.~\ref{app:registries}) forbids redefining a published format. A type admits whichever form it intends its holders to be able to compress; admitting both is permitted and lets a wallet fall back when it cannot produce evidence.

\subsubsection{Reason Body}

CBOR array under tag $\mathsf{SPLITREF\_TAG}$:

\begin{longtable}{clp{8.6cm}}
\textbf{Index} & \textbf{Type} & \textbf{Description} \\
\hline
\endhead
0 & unsigned integer & format version; this specification defines version $1$ \\
1 & byte string & source token identifier $\mathsf{id}_\mathsf{src}$ \\
2 & Certified Transfer & the source's terminal burn transaction $T_\mathsf{burn}^\mathsf{cert}$, in full \\
3 & array & RSMST inclusion proofs, in canonical output-asset order, exactly as in 39044 (Appendix~\ref{app:split-proof-encoding}) \\
\end{longtable}

\noindent The burn is carried in the reason rather than in the evidence set because it is not part of the source token's certificate: a certificate stops immediately before the terminal burn (Sec.~\ref{sec:compression-scope}), and the burn's auxiliary data --- the split manifest $b_\mathcal{M}$, from which the verifier reads the per-asset allocation roots --- must be presented byte-exactly.

Everything the reason must fix is still committed by it: the source identifier, the exact burn, and this output's allocation proofs. The output commitment $d_j$ (Appendix~\ref{app:split-output-commitment}) is unchanged and already binds $\mathsf{id}_\mathsf{src}$.

\subsubsection{Verification}

As $V_{\mathsf{SPLITREF\_TAG}}$ in Alg.~\ref{alg:mint-auth}. The checks are those of Sec.~\ref{sec:split-verification} with check 3 --- recursive verification of the embedded source --- replaced by evidence resolution:

\begin{enumerate}
	\item Decode the reason under the strict rules of Appendix~\ref{app:split-proof-encoding}; ensure $\mathsf{ver} = 1$.
	\item Resolve $\mathsf{id}_\mathsf{src}$ in the evidence set $\mathcal{E}$ (Sec.~\ref{sec:token-evidence}). A missing entry rejects. If the evidence is a complete token, hand it to $\mathsf{collect}$ for worklist verification (Sec.~\ref{sec:nested-verification}) and take $(\mathsf{ty}_\mathsf{src}, \mathcal{A}_\mathsf{src})$ from its genesis and $(h,e)$ from its final pre-burn state; if it is a certificate $\mathcal{K}$, ensure $\Call{VerifyTokenCertificate}{\mathcal{K}, \mathcal{T}, \Gamma} \ne \bot$ and take the same four values from its result. In both cases ensure the resolved identifier equals $\mathsf{id}_\mathsf{src}$ and the network identifier is $\mathcal{T}.\alpha$.
	\item Verify $T_\mathsf{burn}^\mathsf{cert}$ as an ordinary certified transfer (Sec.~\ref{sec:transfer-tx}) out of the state $(\decpred(e), h)$; let $b_\mathcal{M}$ be its auxiliary data. Ensure its recipient predicate is exactly $\predi_\mathsf{burn}(\mathsf{SHA\text{-}256}(b_\mathcal{M}))$.
	\item Apply checks 2 and 5--10 of Sec.~\ref{sec:split-verification} unchanged, reading the source's asset collection from $\mathcal{A}_\mathsf{src}$, the source's type from $\mathsf{ty}_\mathsf{src}$, and the allocation roots from $b_\mathcal{M}$. In particular the output token type MUST be byte-identical to $\mathsf{ty}_\mathsf{src}$ (check 6) and the reconstructed root sum MUST equal the source amount for every asset (check 10).
\end{enumerate}

\paragraph{Security notes.}
\begin{description}
	\item[Same guarantees as 39044] Every check that establishes authenticity, single-spend and value conservation is retained verbatim; only the mechanism by which the verifier obtains the source token changes, and both mechanisms establish the same facts about the same token.
	\item[Identifier binding] $\mathsf{id} = H(\mathsf{salt}, \alpha)$ is unique network-wide and at most one certified genesis exists per identifier (Sec.~\ref{sec:request-validation}), so naming the source by identifier is as tight as embedding it. Substituted evidence for a different token fails the identifier check; evidence for the same token in either form yields the same values.
	\item[Evidence is untrusted] $\mathcal{E}$ is neither certified nor committed. It is checked in full on every verification, and a verifier that cannot verify the evidence it is given rejects, rather than proceeding on partial provenance.
	\item[Amortization] One certificate for the source serves all $k$ outputs of the split and, transitively, their descendants; it is prepared before the split and requires no proving on the transaction path (Sec.~\ref{sec:certificate-preparation}).
\end{description}

\subsection{Merge Reason}\label{app:reason-merge}

\index{token!merging}\index{mint reason!merge}
Merging is the dual of splitting (Sec.~\ref{sec:token-splitting}): $k \ge 2$ source tokens of the same fungible type are burned and one output token carrying the summed assets is minted. It serves consolidation of change and is the aggregation substrate for pooled assets (Sec.~\ref{app:uts-pool}). A type admits merging by including $\mathsf{MERGE\_TAG}$ in $\mathsf{RT}$.

Unlike a split, a merge needs no allocation tree: value conservation is direct arithmetic over the sources' verified payloads. The format is defined in \emph{reference form} (Sec.~\ref{sec:token-evidence}) from the outset --- the reason names its sources by identifier and their evidence travels beside the token --- so a merged token's provenance is compressible and repeated merging does not compound (see \emph{Provenance growth} below).

\subsubsection{Data Structures}

Define $\mathsf{MERGE\_OUTPUT}$ as the ASCII byte string \texttt{"UNICITY\_MERGE\_OUTPUT"}. Let the source tokens have pairwise distinct identifiers, listed in strictly increasing bytewise order: $\mathsf{id}_1 < \cdots < \mathsf{id}_k$, with $2 \le k \le K_\mathsf{max}$, $K_\mathsf{max} = 64$ (in any case bounded by the budget $\Lambda$). For output mint data $D_\mathsf{mint} = (\alpha, \predi', \mathsf{salt}, \mathsf{ty}, \mathsf{reason}, \auxd')$ the \emph{merge output commitment} is
\[
d := \mathsf{SHA\text{-}256}\bigl(\mathsf{CBOR}(
	\mathsf{MERGE\_OUTPUT},\;
	\langle\mathsf{id}_1,\ldots,\mathsf{id}_k\rangle,\;
	\alpha,\;
	\encpred(\predi'),\;
	\mathsf{salt},\;
	\mathsf{id},\;
	\mathsf{ty},\;
	\auxd'
)\bigr),
\]
with term encodings as in Appendix~\ref{app:split-output-commitment} and the source list as a CBOR array of byte strings. The commitment binds the exact source set and the complete output (except the mint reason, excluded to avoid circularity).

The \emph{merge manifest} is the byte string $d$ under semantic tag $\mathsf{MERGE\_MANIFEST\_TAG}$; let $b_\mathcal{M}$ be its canonical encoding. The distinct manifest tag separates merge burns from split burns: a burn commits unambiguously to one protocol.

The mint reason is a CBOR array under tag $\mathsf{MERGE\_TAG}$: the format version ($1$), the source identifiers $\langle \mathsf{id}_1, \ldots, \mathsf{id}_k\rangle$ in strictly increasing order, and the $k$ certified terminal burn transactions in the same order. The sources themselves are resolved from the evidence set (Sec.~\ref{sec:token-evidence}), each as a complete token or as a token certificate.

\subsubsection{Merge Protocol}

\begin{enumerate}
	\item Verify every source token completely, including payloads and mint reasons. All sources MUST have byte-identical token type $\mathsf{ty}$ and network identifier $\alpha$; merging is not a type-conversion mechanism.
	\item Compute the output collection: for every asset identifier occurring in any source, $v_\mathsf{out}(\mathsf{aid}) = \sum_i v_i(\mathsf{aid})$ (missing entries are zero), rejecting any sum $\ge 2^{256}$. The output payload $\auxd'$ is exactly this canonical summed collection --- no asset omitted, none introduced.
	\item Prepare the output mint fields, derive $\mathsf{id} = H(\mathsf{salt}, \alpha)$, compute $d$ and $b_\mathcal{M}$. Durably store the output mint fields and manifest before proceeding.
	\item Burn each source with auxiliary data exactly $b_\mathcal{M}$ and recipient predicate $\predi_\mathsf{burn}(\mathsf{SHA\text{-}256}(b_\mathcal{M}))$ (Sec.~\ref{sec:burn-transaction}). Burns may execute in any order; an interruption strands no value permanently, since the pre-committed output can be minted whenever all burns eventually complete.
	\item Mint the output with the reason naming the burned sources and carrying their burn transactions, and assemble the evidence set for those sources (Sec.~\ref{sec:token-evidence}).
\end{enumerate}

\subsubsection{Merge Verification}

As $V_{\mathsf{MERGE\_TAG}}$ in Alg.~\ref{alg:mint-auth}; every failed check rejects.
\begin{enumerate}
	\item Decode the reason under the strict rules of Appendix~\ref{app:split-proof-encoding}; require $2 \le k \le K_\mathsf{max}$.
	\item Decode the output payload: $\mathcal{A}_\mathsf{out} \gets \Call{Assets}{D_\mathsf{mint}.\mathsf{ty}, D_\mathsf{mint}.\auxd'}$, a non-empty canonical collection of positive amounts.
	\item Require the identifiers $\mathsf{id}_i$ strictly increasing in $i$. For each: resolve $\mathsf{id}_i$ in the evidence set $\mathcal{E}$ (Sec.~\ref{sec:token-evidence}) --- a missing entry rejects --- either handing a complete token to $\mathsf{collect}$ for worklist verification (Sec.~\ref{sec:nested-verification}) or verifying a token certificate by $\Call{VerifyTokenCertificate}{}$, which must not return $\bot$ (Sec.~\ref{sec:certificate-verification}). Take $(\mathsf{ty}_i, \mathcal{A}_i, h_i, e_i)$ from the resolved evidence; require the resolved identifier to equal $\mathsf{id}_i$, the type to be byte-identical to $D_\mathsf{mint}.\mathsf{ty}$, and the network identifier to be $\mathcal{T}.\alpha$.
	\item Recompute $d$ from $\langle\mathsf{id}_1,\ldots,\mathsf{id}_k\rangle$ and the output mint fields other than $\mathsf{reason}$; form $b_\mathcal{M}$. Verify each carried burn transaction as an ordinary certified transfer out of $(\decpred(e_i), h_i)$ (Sec.~\ref{sec:transfer-tx}), and require its auxiliary data to be exactly $b_\mathcal{M}$ and its recipient predicate exactly $\predi_\mathsf{burn}(\mathsf{SHA\text{-}256}(b_\mathcal{M}))$.
	\item Conservation: compute the canonical summed collection over $\mathcal{A}_1, \ldots, \mathcal{A}_k$, rejecting overflow; require it to equal $\mathcal{A}_\mathsf{out}$ exactly --- the same asset identifiers with the same amounts.
\end{enumerate}

\paragraph{Security notes.}
\begin{description}
	\item[Authenticity] Every source is verified under the same $\Gamma$, $\mathcal{T}$ and remaining budget as the output (Sec.~\ref{sec:verifier-registry}), whether it is supplied in full or as a certificate.
	\item[No double counting] Strictly increasing identifiers exclude duplicate sources within a merge; the single-spend property (Sec.~\ref{sec:request-validation}) excludes burning a source twice at all; and $d$ binds each burn to this exact source set and output, so a burn cannot be replayed in a different merge.
	\item[No inflation] Exact conservation with overflow rejection; byte-identical type excludes conversion (as in split verification, Sec.~\ref{sec:split-verification}, check~6).
	\item[Fixed ownership] The output recipient, salt and payload are committed by $d$ before any irreversible burn.
	\item[Value risk, not inflation risk] As with splits, failing to mint the pre-committed output destroys value but can never create or redirect it; the merging party is responsible for durable storage of the pre-committed data.
\end{description}

\paragraph{Provenance growth.} With sources supplied in full, a merge carries $k$ complete tokens and repeated merging grows provenance multiplicatively, bounded by $\Lambda$. The intended composition is with compression (Secs.~\ref{sec:history-compression}, \ref{sec:provenance-compression}): sources SHOULD be supplied as token certificates, prepared in advance by each source holder. Each source then costs one constant-size certificate irrespective of its own history and ancestry, and the merged token's own certificate --- prepared later --- folds all $k$ of them into a single proof. What survives is the issuance roots: each independent lineage entering the merge contributes one genesis that every verifier checks in the clear (Sec.~\ref{sec:token-certificate}), so a merged token's cost is proportional to the number of distinct issuance events behind it rather than to its transaction history. Each source must in any case present its terminal burn in the clear.

\subsection{Conversion Reason}\label{app:reason-convert}

\index{mint reason!conversion}
Split and merge conserve value \emph{within} one token type. The conversion format carries value \emph{across} a type boundary: a mint of type $\mathsf{ty}_\mathsf{out}$ is authorized by a certified transition (transfer or burn) of a \emph{source} token of a different type $\mathsf{ty}_\mathsf{in}$, whose own validation rules committed to the output and debited whatever the output credits. The canonical use is pool entry and exit --- share tokens minted against pool deposits, plain tokens minted against pool withdrawals (Appendix~\ref{app:autonomous-pools}).

The pattern is the familiar one: the authorizing transition's auxiliary data contains, as mandated by $\mathsf{ty}_\mathsf{in}$'s type rules, a \emph{conversion output commitment}
\[
d := \mathsf{SHA\text{-}256}\bigl(\mathsf{CBOR}(
	\mathsf{CONVERT\_OUTPUT},\;
	\mathsf{id}_\mathsf{src},\;
	\alpha,\;
	\encpred(\predi'),\;
	\mathsf{salt},\;
	\mathsf{id},\;
	\mathsf{ty}_\mathsf{out},\;
	\auxd'
)\bigr)
\]
over the complete output mint fields except the reason ($\mathsf{CONVERT\_OUTPUT} = \texttt{"UNICITY\_CONVERT\_OUTPUT"}$; term encodings as in Appendix~\ref{app:split-output-commitment}). One commitment authorizes exactly one mint: $d$ fixes the salt, hence the output identifier, hence a genesis state certified at most once.

\begin{itemize}
	\item \textbf{Type opt-in, both sides}: $\mathsf{ty}_\mathsf{out}$'s manifest names its admissible \emph{conversion sources} $\mathcal{CS}$ --- the set of source types (or a rule identifying them) whose transitions may mint $\mathsf{ty}_\mathsf{out}$; $\mathsf{ty}_\mathsf{in}$'s transfer rules define which of its transitions commit conversion outputs, where in the transition record the commitment sits, and what the transition debits in exchange. Verifying a conversion therefore requires \emph{both} type definitions in $\Gamma.\mathcal{TD}$; either one missing is a fail-shut rejection (Sec.~\ref{sec:mint-authorization}).
	\item \textbf{Reason body}: CBOR array under tag $\mathsf{CONVERT\_TAG}$: $[\mathsf{ver}, i, \mathcal{L}_\mathsf{src}]$ --- format version ($1$), the index $i \ge 1$ of the committing transition, and the complete source token up to and including transition $i$.
\end{itemize}

\textbf{Verification} (as $V_{\mathsf{CONVERT\_TAG}}$ in Alg.~\ref{alg:mint-auth}):
\begin{enumerate}
	\item Decode $(\mathsf{ver}, i, \mathcal{L}_\mathsf{src})$; ensure $\mathsf{ver} = 1$.
	\item Let $\mathsf{ty}_\mathsf{in}$ be the source's genesis token type; ensure $\mathsf{ty}_\mathsf{in} \in \mathcal{D}_{\mathsf{ty}_\mathsf{out}}.\mathcal{CS}$ and the source's network identifier equals $\mathcal{T}.\alpha$; hand $\mathcal{L}_\mathsf{src}$ to $\mathsf{collect}$ for worklist verification (which re-runs $\mathsf{ty}_\mathsf{in}$'s transition rules over the source's entire history).
	\item Recompute $d$ from $\mathsf{id}(\mathcal{L}_\mathsf{src})$ and the output mint fields other than $\mathsf{reason}$; ensure transition $i$ of $\mathcal{L}_\mathsf{src}$ exists and its record contains exactly $d$ at the position mandated by $\mathsf{ty}_\mathsf{in}$'s rules.
\end{enumerate}

\paragraph{Soundness condition.} The format itself is conservation-neutral: it proves only that a certified source transition committed to this exact output. \emph{What} the output is worth, and that the source debited it, are properties of $\mathsf{ty}_\mathsf{in}$'s transition rules --- which the worklist verification of $\mathcal{L}_\mathsf{src}$ re-checks in full. A source type whose rules commit outputs without debiting them simply defines an inflationary conversion; $\mathsf{ty}_\mathsf{out}$'s registration of $\mathcal{CS}$ is precisely the statement that its verifiers accept those source rules as sound. (Split and merge are, in retrospect, same-type special cases of this pattern with burn-bound rather than transition-bound commitments; they remain separately specified.)

\subsection{Proof-Carrying Reason}\label{app:reason-proof}

\index{mint reason!zero-knowledge proof}
Minting is gated by a succinct (optionally zero-knowledge) proof of an arbitrary, type-defined statement: the general form of which the bridge backing proof and the split proof are hand-rolled special cases. The token type is matched with a verification key: the type definition pins the proof system $\mathsf{PS}$, the verification key hash $h_\mathsf{vk}$, and the \emph{public-input schema} --- the deterministic derivation of the public input vector from the genesis and the reason body.

\begin{itemize}
	\item \textbf{Type parameters}: $(\mathsf{PS}, h_\mathsf{vk}, \text{public-input schema})$, all frozen in the definition manifest. Content-addressed derivation $\mathsf{ty} = H(\mathsf{TYPE\_DEF}, b_\mathcal{D})$ is RECOMMENDED, making the token type a commitment to its own verification key.
	\item \textbf{Reason body}: CBOR array under tag $\mathsf{PROOF\_TAG}$:
\end{itemize}

\begin{longtable}{clp{8.6cm}}
\textbf{Index} & \textbf{Type} & \textbf{Description} \\
\hline
0 & unsigned integer & format version; this specification defines version $1$ \\
1 & byte string & verification key $\mathsf{vk}$ \\
2 & CBOR structure & auxiliary public inputs $x_\mathsf{aux}$, per the type's schema \\
3 & byte string & proof $\pi_\mathsf{zk}$ \\
\end{longtable}

\textbf{Verification} (as $V_{\mathsf{PROOF\_TAG}}$ in Alg.~\ref{alg:mint-auth}):
\begin{enumerate}
	\item Decode $(\mathsf{ver}, \mathsf{vk}, x_\mathsf{aux}, \pi_\mathsf{zk})$; ensure $\mathsf{ver} = 1$ and $H(\mathsf{vk}) = \mathcal{D}_\mathsf{ty}.h_\mathsf{vk}$.
	\item $x \gets \Call{PublicInputs}{m_\mathsf{mint}(D_\mathsf{mint}), x_\mathsf{aux}}$ per the type's schema. The schema MUST include the genesis commitment $m_\mathsf{mint}$ (Sec.~\ref{app:genesis-commitment}) in $x$, so the proven statement binds this exact mint.
	\item Ensure $\mathsf{PS}.\mathsf{Verify}(\mathsf{vk}, x, \pi_\mathsf{zk}) = 1$.
\end{enumerate}

\paragraph{Security notes.} Soundness reduces to the proof system's knowledge soundness and the immutability of $h_\mathsf{vk}$; the upgrade discipline is that of Sec.~\ref{sec:bridge-security}: a relation is never swapped in place, a new relation is a new type. The semantics of the proven statement is part of the frozen definition --- a recipient who cannot obtain and understand the statement's specification MUST NOT register the type (fail-shut). Proof verification cost is constant for SNARKs and is charged against the budget $\Lambda$.

\section{Issuance Policies}\label{app:issuance-policies}

\index{issuance policy}\index{issuance!direct}\index{issuance!derived}
A certified genesis is either a \emph{direct issuance} or a \emph{derived issuance}. A direct issuance creates an item or an amount of value that no certified token of the type held before. A derived issuance re-mints value that certified tokens already hold: a split, reference-form split or merge of tokens of the same type, or a conversion from a declared source type (Sec.~\ref{app:reason-convert}). The conservation checks of its reason format bound a derived issuance by its sources. The admissible reason set of a type accordingly decomposes as
\[
\mathsf{RT} = \mathsf{RT}_\mathsf{iss} \cup \mathsf{RT}_\mathsf{der},
\qquad
\mathsf{RT}_\mathsf{iss} \cap \mathsf{RT}_\mathsf{der} = \emptyset .
\]
The usage standard of the type (Sec.~\ref{app:standard-types}) fixes $\mathsf{RT}_\mathsf{der}$. The \emph{issuance policy} fixes $\mathsf{RT}_\mathsf{iss}$ together with the checks that $\Call{VerifyIssuance}{}$ applies to direct geneses. The two choices are independent. The standard determines what a token is and which operations preserve it; the policy determines who may create tokens and how many can exist. A deployed type freezes one standard, one policy and the parameters of both in its manifest (Sec.~\ref{app:type-registry}).

\begin{longtable}{>{\raggedright\arraybackslash}p{1.9cm}>{\raggedright\arraybackslash}p{2.8cm}>{\raggedright\arraybackslash}p{2.2cm}>{\raggedright\arraybackslash}p{3.0cm}>{\raggedright\arraybackslash}p{3.4cm}}
\textbf{Policy} & $\mathsf{RT}_\mathsf{iss}$ & \textbf{Parameters} & \textbf{Bound on direct issuance} & \textbf{Ethereum practice} \\
\hline
\endhead
Open & $\{\bot\}$ & none & none & public mint function \\
Authorized & $\{\mathsf{ISSUER\_TAG}\}$ & $\predi_\mathsf{issue}$ & at the issuer's discretion & owner- or role-restricted mint \\
Capped & one tag of the selected profile & profile parameters, $\predi_\mathsf{issue}$ & at most $N$ items or $X$ units per network & capped supply (e.g.\ \texttt{ERC20Capped}, fixed-size collections) \\
Bridged & $\{39048\}$ & $h_\mathsf{cfg}$ & value locked on the external chain & lock-and-mint bridge \\
Proof-gated & $\{\mathsf{PROOF\_TAG}\}$ & $\mathsf{PS}$, $h_\mathsf{vk}$, schema & as defined by the proven statement & mint guarded by a verifier contract \\
\end{longtable}

On Ethereum these restrictions are code in the token contract and are enforced once, when the minting transaction executes. Here every recipient enforces them again, by verifying the genesis of each token it accepts against the frozen definition (Alg.~\ref{alg:mint-auth}). No party keeps a supply counter, so a bound on direct issuance is a property that each token proves about itself and its type, not a figure that anyone can query (Sec.~\ref{app:capped-availability}).

\subsection{Open Issuance}\label{app:policy-open}

$\mathsf{RT}_\mathsf{iss} = \{\bot\}$, and $\Call{VerifyIssuance}{}$ imposes no condition on direct geneses beyond the rules of the standard. Anyone can mint any payload that $\Call{ParseAux}{}$ accepts. A token of an open type proves that its identifier was fresh and that its ownership chain is certified; it proves nothing about who minted it. Byte-identical payloads can be minted any number of times under distinct identifiers, so authenticity claims must be carried inside the payload and assessed by the application, for example the creator attribution of UTS-2 (Sec.~\ref{app:nft-attribution}). Where two tokens carry the same content, the certification times $\mathsf{time}(\pi_0)$ of their geneses establish which was minted first.

Open issuance of a fungible asset allows anyone to inflate its supply without limit. Such a type has no value semantics, and wallets MUST NOT present its tokens as carrying value; it is admissible only for testing (e.g.\ a test-network faucet).

\subsection{Authorized Issuance}\label{app:policy-authorized}

$\mathsf{RT}_\mathsf{iss} = \{\mathsf{ISSUER\_TAG}\}$. The manifest pins the issuance predicate $\predi_\mathsf{issue}$, and each direct genesis carries the issuer's unlocking argument over its genesis commitment (Sec.~\ref{app:reason-issuer}). The standard may add per-mint payload constraints, e.g.\ the permitted asset identifiers of a fungible type. The number and amount of direct issuances are unbounded: recipients rely on the issuer's discipline, as holders of an Ethereum token rely on the owner of a restricted mint function. Because $\predi_\mathsf{issue}$ is an arbitrary predicate, the issuer may be a single key, a $k$-of-$n$ committee or a time-limited authorization. A compromised issuance key cannot be revoked retroactively (Sec.~\ref{app:reason-issuer}).

\subsection{Bridged Issuance}\label{app:policy-bridged}

$\mathsf{RT}_\mathsf{iss} = \{39048\}$ with the bridge configuration hash $h_\mathsf{cfg}$ as parameter. Each direct genesis proves a lock of the same value on an external chain (Sec.~\ref{sec:bridged-assets}), so outstanding value is bounded by the locked value under the bridge's security assumptions. The bridge specification covers fungible assets only; bridging of non-fungible items is not specified.

\subsection{Proof-Gated Issuance}\label{app:policy-proof}

$\mathsf{RT}_\mathsf{iss} = \{\mathsf{PROOF\_TAG}\}$. Each direct genesis carries a succinct proof of a type-defined statement over its genesis commitment (Sec.~\ref{app:reason-proof}). Any bound on issuance is a property of that statement, and a recipient who registers the type accepts the statement's specification as part of the frozen definition.

\subsection{Capped Issuance}\label{app:uts-capped}

\index{token type!capped issuance}\index{issuance limit}
The capped policy bounds direct issuance of fungible (UTS-1) and non-fungible (UTS-2) types. Every valid direct issuance consumes one of a bounded set of genesis certification states; fungible descendants additionally conserve their source value. Recipients verify these facts locally from the frozen definition and token evidence, without observing every issuance. A cap applies separately on each network $\alpha$; the same definition on several networks does not establish a combined cross-network cap.

\subsubsection{Profile Selection}\label{app:capped-selection}

The manifest MUST select exactly one of the following issuance profiles and its parameters. \textbf{Profile F is the default for fungible assets; profile V is the default for NFT collections.} The choice fixes the reason formats and validation rules for the lifetime of the type.

\begin{longtable}{>{\raggedright\arraybackslash}p{2.6cm}>{\raggedright\arraybackslash}p{4.9cm}>{\raggedright\arraybackslash}p{6.6cm}}
\textbf{Profile} & \textbf{Pick when} & \textbf{Construction and cost} \\
\hline
\endhead
F: Singleton genesis & Fixed fungible supply; distribute a treasury over time. & One initial token of value $X$, then split/merge. Constant-size issuance commitment; descendants share one issuance root. \\
V: VRF serials & NFTs minted as needed, or independent fungible mints of equal denomination, without a slot table. & Frozen range and VRF key; one VRF proof per direct mint. Requires the specified VRF verifier and confidential key custody. \\
S: Committed slots & Arbitrary pre-committed denominations or NFT contents; or a deployment needing hash-based commitments instead of a VRF. & Generate $N$ secret salts and a tree; carry a membership or sum proof per mint. Full table distribution is optional for cap verification. \\
P: Public serials & Publicly enumerable NFT mint identities are required and permanent blocking is acceptable. & Frozen range and public hash derivation; no VRF or slot proof. Not the default for issuer-controlled issuance. \\
\end{longtable}

Independent mints (V and S) avoid spending a shared treasury state and its accumulated split ancestry. F can obtain parallel distribution by splitting the treasury into allocations beforehand; its full supply is then already issued. S adds arbitrary per-slot commitments and permits separate custody of individual salts. V fixes serial membership without an $O(N)$ commitment table. If the required VRF implementation is unavailable, choose S for NFTs; do not silently substitute P or an unchecked VRF.

\subsubsection{Common Rules}

\begin{itemize}
	\item The type identifier MUST be content-addressed (Sec.~\ref{sec:token-type}), committing to the profile, cap, admissible reasons and all issuance parameters. Cap values belong to the consensus-relevant manifest; display metadata is not evidence of a cap.
	\item A fungible deployment pins one $\mathsf{aid}$ and uses the canonical asset payload rules of UTS-1. Its $\mathsf{RT}$ contains only its selected direct-issuance tag and a subset of $\{39044,\mathsf{SPLITREF\_TAG},\mathsf{MERGE\_TAG}\}$. It SHOULD admit and prefer the reference split form for compressible provenance. Profile F MUST admit at least one split format.
	\item An NFT deployment uses the payload rules of UTS-2 and admits only its selected direct-issuance tag: no split, merge or conversion. Its serial is the verified reason's $j$, immutable as part of the genesis; it is not a field of the UTS-2 payload (Sec.~\ref{app:nft-collection}). Distinct serials do not by themselves prove distinct content.
	\item Absent reasons and all other issuance paths are excluded. In particular, V and S MUST NOT also admit unrestricted issuer-signature mints; F admits that format only with its singleton restrictions below. Combining profiles requires a separately specified definition with a proof of the combined bound.
	\item Compression is a representation change, not an issuance format. The original direct-issuance geneses remain the issuance roots presented in the clear and checked by recipients (Sec.~\ref{sec:provenance-compression}). A recipient of a profile-F descendant therefore sees and verifies the original singleton mint even when its split ancestry is compressed.
\end{itemize}

For fungible assets, $X$ bounds the value created by direct issuance and consequently all outstanding value, including an issuer's undistributed treasury. It does not bound the number of containers or the cumulative value of repeated split/merge mints. A release from a preminted treasury changes distribution, not total supply. Permanent burns can lower outstanding value; the cap does not guarantee that $X$ is ever issued or remains spendable.

\subsubsection{Profile F: Singleton Genesis (Fungible Default)}\label{app:uts-singleton}

The manifest pins $\mathsf{aid}$, $1 \le X < 2^{256}$, $\predi_\mathsf{issue}$ and one salt commitment $h_\mathsf{single}$. Define $\mathsf{SINGLETON\_SALT}$ as the ASCII byte string \texttt{"UNICITY\_SINGLETON\_SALT"}. Deployment and minting proceed as follows:
\begin{enumerate}
	\item Sample a 32-byte salt $s$ with at least 128 bits of min-entropy; keep it confidential. Compute
\[
h_\mathsf{single}:=\mathsf{SHA\text{-}256}(\mathsf{CBOR}(\mathsf{SINGLETON\_SALT},s)),
\]
where the input is one deterministic CBOR array of byte strings. Freeze the manifest and derive $\mathsf{ty}$. The commitment omits the future mint transaction and $\mathsf{ty}$, avoiding a circular definition.
	\item Construct the single direct mint with salt $s$, payload exactly the canonical encoding of $\{(\mathsf{aid},X)\}$, and an issuer-signature reason (Sec.~\ref{app:reason-issuer}) authorizing its complete genesis commitment and chosen first owner. Obtain certification before disclosing the salt or mint data.
	\item Distribute amounts by ordinary transfers and value-conserving splits; merge to consolidate where admitted. The same split rules bind every output's type and allocation to the consumed source.
\end{enumerate}

For every $\mathsf{ISSUER\_TAG}$ genesis, $\Call{VerifyIssuance}{}$ MUST recompute the salt commitment and require equality to $h_\mathsf{single}$, require the exact pinned asset and amount $X$, and enforce the pinned issuance predicate through the reason verifier. These singleton checks apply to \emph{direct issuance only}; split and merge outputs are checked by their conservation verifiers and may have different salts and amounts.

Under the hash assumptions, a fixed salt commitment permits one salt, hence one genesis certification state per network. The Service's uniqueness rule permits at most one direct mint even if the issuer signs alternatives. Every admitted descendant conserves that initial value, establishing the cap without a table, supply query or issuer promise. Merely minting a large UTS-1 token without these additional restrictions would not enforce a cap. All $X$ units are issued at genesis; protecting an undistributed balance is ordinary ownership-key custody, and a timed release policy requires additional spending rules.

\subsubsection{Profile V: VRF Serials (NFT Default)}\label{app:uts-vrf}

Use the parameters, derivation and reason verification of Sec.~\ref{app:reason-vrf}, with $\mathsf{VRF\_TAG}$ as the sole direct-issuance tag. Generate an independent VRF key pair, freeze its public key and the serial range in the manifest, and compute each proof and salt only when needed (or retain precomputed values confidentially). Authorize and certify that serial's mint before delivering its data and proof. Each NFT proves that its serial belongs to the range and has only one permitted mint identity; its content is chosen at mint time under the issuer's authorization. Pick S when arbitrary content must instead be committed before issuance.

For an independently issued fungible series, the manifest additionally pins a positive denomination $v$ and $X=Nv<2^{256}$. Type registration MUST check this equality without overflow, and $\Call{VerifyIssuance}{}$ MUST require every direct mint's payload to be exactly $\{(\mathsf{aid},v)\}$. At most $N$ direct mints can then issue at most $X$ units; later splits and merges conserve value. Choose this variant when parallel or staged direct issuance matters and a uniform denomination is sufficient; choose F for ordinary treasury distribution or S for arbitrary denominations. Free choice of the amount at mint time is not permitted by this profile.

\subsubsection{Profile S: Pre-Committed Slots}\label{app:uts-slots}

Use $\mathsf{SLOT\_TAG}$ as the sole direct-issuance tag and construct the commitments of Sec.~\ref{app:reason-slot}. Generate and back up the secret salts, fix the per-slot constraints, build the tree and freeze its root and mode in the manifest. Retain the tree data needed to construct proofs. For each mint, open one commitment with its salt and proof, obtain issuer authorization and certification, and only then disclose the token. A recipient verifies its own proof; it does not need the rest of the table.

Fungible deployments MUST use value mode and pin $\mathsf{aid}$ and $X$: each proof authenticates both the slot's exact denomination and a reconstructed total equal to $X$. The full audit table is optional even for heterogeneous denominations. NFT deployments MUST use count mode, with serials $1,\ldots,N$ and either of the following constraints per slot:
\begin{description}
	\item[Pre-committed content] $b_j$ fixes the exact UTS-2 payload bytes, e.g.\ a metadata item with inline media or a document link with its content digest (Sec.~\ref{app:nft-items}). A disclosed table lets readers audit the entire collection before issuance. Commitment alone does not establish that content hashes are pairwise distinct: any claim of distinct committed contents requires checking the complete authenticated table for duplicates.
	\item[Open content] $b_j=\bot$ leaves content to mint time while fixing the serial and mint identity. Choose this when hash-based slot proofs or individual salt custody are preferred to the VRF profile.
\end{description}

\subsubsection{Profile P: Public NFT Serials}\label{app:uts-public-serial}

Use $\mathsf{SERIAL\_TAG}$ as the only admissible reason and the range and derivation of Sec.~\ref{app:reason-serial}. Freeze the definition, derive the public mint salt for a serial, then authorize and certify its mint. This is the least complex serial construction, but any observer can derive every unused certification state as soon as the definition is available. Select it only when that public enumerability is intentional and the application accepts that some or all serials may be permanently blocked. A trivially-true issuance predicate permits public claims, with the same limitation; it does not guarantee a valid token for every serial.

\subsubsection{Availability and Issuance Audits}\label{app:capped-availability}

\paragraph{Predictable mint states can be blocked.} The minting private key is publicly derivable from the token identifier (Appendix~\ref{app:el-constants}). An attacker who learns an unused identifier can submit a signed request for its genesis state with an arbitrary transaction hash. The Service checks the request predicate and state uniqueness, not the token type's mint reason or payload (Sec.~\ref{sec:request-validation}). Such a request can therefore consume the state forever without creating a token any recipient accepts. An issuer signature in the reason does not prevent this denial of issuance. It does prevent acceptance of unauthorized tokens; the cap still holds.

F and S hide unused identities behind high-entropy salt commitments; V derives them using a confidential VRF key. For these protections to apply, the issuer MUST obtain mint certification before exposing the salt, identifier, VRF output or proof to untrusted parties, including an intended recipient. Certification requests expose only the required hash images and mint public key, not those preimages. Disclosure of the completed token is safe because its genesis state is already consumed. Secret loss or premature disclosure can destroy availability but cannot increase the cap. P deliberately gives up this protection. Fixing public-state blocking at the Service would require a separately specified authorization mechanism; it is not assumed here.

\paragraph{A cap is not a public issuance history.} A frozen definition and local proofs establish an upper bound without a shared public blockchain. A disclosed mint's certificate establishes its certification time, but unused-looking private slots or serials may already have been minted without public disclosure. Neither an audit table nor a VRF proves the complete current issued count, circulating supply, or absence of undisclosed mints.

Applications requiring an enforced issuance schedule MUST freeze deterministic serial-specific time windows or release rules in $\Call{VerifyIssuance}{}$ and check each direct mint against $\mathsf{time}(\pi_0)$. These checks must apply independently of issuer discretion; an ordinary signature alone is not a schedule. V and S support such checks without sharing a spendable treasury state. F instead separates the time of initial issuance from later treasury distribution. A publicly auditable \emph{complete} history additionally needs a publication and completeness mechanism, which these profiles do not specify. Do not infer such a history from the cap or the timestamps of a subset of tokens.

\section{Standard Token Types}\label{app:standard-types}

\index{token type!standard}
Standard token types are identified as \emph{UTS-$n$} (Unicity Token Standard); numbers are assigned in the registry. A standard is defined by \emph{usage}. It fixes the payload decoder $\Call{ParseAux}{}$, the derived-issuance formats $\mathsf{RT}_\mathsf{der}$ that the usage requires, any per-transfer checks, and the metadata conventions. It leaves the issuance policy open (Sec.~\ref{app:issuance-policies}), restricted only to the policies that are meaningful for the usage. A deployed type instantiates one standard and one policy with concrete parameters and has its own type identifier. We write, e.g., UTS-1/capped-F for a fungible type under the singleton profile of capped issuance.

The standards aim at functional parity with the corresponding Ethereum token standards, so that their concepts and field names are familiar. They are not interface-compatible. An Ethereum token is an entry in the ledger of a contract; a Unicity token is a bearer object that its recipient verifies. Operations that read or modify a shared ledger (balance queries, total supply, allowances) have no direct counterpart. The correspondence tables below give the substitute where one exists.

\begin{longtable}{>{\raggedright\arraybackslash}p{1.1cm}>{\raggedright\arraybackslash}p{1.9cm}>{\raggedright\arraybackslash}p{2.3cm}>{\raggedright\arraybackslash}p{2.3cm}>{\raggedright\arraybackslash}p{2.3cm}>{\raggedright\arraybackslash}p{3.0cm}}
\textbf{Std.} & \textbf{Usage} & \textbf{Ethereum counterpart} & $\Call{ParseAux}{}$ & $\mathsf{RT}_\mathsf{der}$ & \textbf{Issuance policies} \\
\hline
\endhead
UTS-0 & open container & none & any byte string & $\emptyset$ & open (canonical instance), authorized \\
UTS-1 & fungible token & ERC-20; fungible part of ERC-1155 & canonical asset collection & split formats, optionally merge and conversion & authorized, capped (F, V, S), bridged, proof-gated \\
UTS-2 & non-fungible token & ERC-721 and its metadata schema & NFT item & $\emptyset$ & open, authorized, capped (V, S, P), proof-gated \\
UTS-3 & pooled assets & liquidity pools and their share tokens & UTS-1 collections; pool state & conversion & see Sec.~\ref{app:uts-pool} \\
\end{longtable}

\subsection{UTS-0: Open Container}\label{app:uts-blank}

A type for arbitrary, possibly absent payload.

\begin{itemize}
	\item $\Call{ParseAux}{}$ accepts any byte string, including an absent payload.
	\item $\mathsf{RT}_\mathsf{der} = \emptyset$.
	\item \textbf{Canonical instance}: open issuance, $\mathsf{RT} = \{\bot\}$ and $\Call{VerifyIssuance}{} \equiv 1$. Its identifier is a single well-known value: the content-addressed identifier (Sec.~\ref{sec:token-type}) of the canonical UTS-0 definition manifest, recorded in the registry. Other instances MAY use authorized issuance; the type then proves which issuer authorized a payload but still assigns the payload no meaning.
\end{itemize}

An open-payload token inherits the platform guarantees: a unique identifier, a certified ownership chain and double-spend protection. It has no issuance semantics beyond those of its policy. It suits application-defined artifacts (tickets, attestations, session tokens, testing) where the application assigns and validates the meaning of $\auxd'$. Generic wallets MUST NOT present UTS-0 tokens as carrying value.

A wallet MAY display an open token whose payload is a recognized UTS-2 item (Sec.~\ref{app:nft-validity}) as an NFT. That interpretation is presentation only. A payload that fails recognition does not make the token invalid, and a recognized payload in a UTS-0 type carries none of the type-level guarantees of UTS-2.

\subsection{UTS-1: Fungible Token}\label{app:uts-fungible}

The fungible token standard, the analogue of ERC-20: a divisible container of amounts of one or more assets.

\begin{itemize}
	\item $\Call{ParseAux}{}$ is the canonical asset-collection decoding $\Call{Assets}{}$ (Sec.~\ref{sec:asset-structure}, canonical rules of Appendix~\ref{app:split-proof-encoding}). The payload rules pin the set $\mathcal{AID}_\mathsf{ty}$ of permitted asset identifiers; a collection with any other identifier rejects.
	\item $\mathsf{RT}_\mathsf{der} \subseteq \{39044, \mathsf{SPLITREF\_TAG}, \mathsf{MERGE\_TAG}, \mathsf{CONVERT\_TAG}\}$ and contains at least one split format. Divisibility is the defining property of the usage: an owner must be able to pay any amount up to the value of a token (Sec.~\ref{sec:token-splitting}). A type SHOULD admit $\mathsf{MERGE\_TAG}$ for consolidating change (Sec.~\ref{app:reason-merge}). A type intended for high circulation SHOULD admit the reference split form $\mathsf{SPLITREF\_TAG}$ (Sec.~\ref{app:reason-splitref}), and its wallets SHOULD prefer it, so that split-minted provenance is compressible. $\mathsf{CONVERT\_TAG}$ is admitted only by types that participate in pools (Sec.~\ref{app:autonomous-pools}).
	\item $\Call{VerifyIssuance}{}$ applies the checks of the issuance policy to direct geneses and requires their asset identifiers to lie in $\mathcal{AID}_\mathsf{ty}$. Derived geneses are constrained by their reason verifiers: the output type equals the source type, and value is conserved per asset identifier (Sec.~\ref{sec:split-verification}).
	\item \textbf{Issuance policies}: authorized issuance for assets issued on demand; capped profile F for a fixed supply distributed from a treasury, profile V for independently minted equal denominations, profile S in value mode for pre-committed denominations; bridged issuance for assets locked on an external chain; proof-gated issuance for type-defined rules. Open issuance is admissible only for testing (Sec.~\ref{app:policy-open}).
	\item \textbf{Metadata}: a registry entry with $\texttt{assetKind} = \texttt{"fungible"}$, carrying name, symbol, decimals and icons (Sec.~\ref{app:asset-metadata}).
\end{itemize}

Amounts are integers in minimal units; $\texttt{decimals}$ scales the display only, as in ERC-20. A type that permits several asset identifiers lets one token carry a vector of balances, one per asset, and split and merge conserve each component separately. This corresponds to the fungible part of ERC-1155, in which one contract manages several fungible token identifiers.

\begin{longtable}{>{\raggedright\arraybackslash}p{4.2cm}>{\raggedright\arraybackslash}p{10.0cm}}
\textbf{ERC-20} & \textbf{UTS-1} \\
\hline
\endhead
\texttt{name}, \texttt{symbol}, \texttt{decimals} & registry metadata of the type (Sec.~\ref{app:asset-metadata}); presentation only \\
\texttt{totalSupply} & no global counter; capped issuance proves an upper bound $X$, not the outstanding amount (Sec.~\ref{app:capped-availability}) \\
\texttt{balanceOf} & sum of $\Call{Assets}{}$ over the tokens a wallet holds, computed by the holder; not observable by others \\
\texttt{transfer} & transfer of a token carrying exactly the amount, if necessary after splitting off the amount and change (Sec.~\ref{sec:token-splitting}) \\
\texttt{approve}, \texttt{allowance}, \texttt{transferFrom} & no allowance ledger; delegated and conditional spending is expressed by the owner predicate: multi-signature, threshold, time lock, hash lock or P2SH (Sec.~\ref{sec:standard-predicates}) \\
\texttt{Transfer} event & certified transaction, visible to the parties that hold the token; nothing is published \\
mint (extension) & direct issuance under the type's policy \\
burn (\texttt{ERC20Burnable}) & burn transaction (Sec.~\ref{sec:burn-transaction}) \\
\end{longtable}

\subsection{UTS-2: Non-Fungible Token}\label{app:uts-nft}

The non-fungible token standard, the analogue of ERC-721. Each token is one unique item. Its identity is the token identifier $\mathsf{id}$ (Sec.~\ref{sec:mint-tx}), which is unique on the network; its description is the genesis payload $\auxd'_0$, which is immutable. The payload format is the Sphere NFT metadata format\footnote{\url{https://github.com/unicity-sphere/sphere-sdk/blob/main/docs/NFT-METADATA.md}; it contains the byte-level grammars and test vectors that this section omits.}, version~1. Its fields follow the ERC-721 metadata JSON schema, encoded as tagged deterministic CBOR arrays.

\begin{itemize}
	\item $\Call{ParseAux}{}$ is recognition of an NFT item together with the collection-identity rule (Secs.~\ref{app:nft-validity}, \ref{app:nft-collection}).
	\item $\mathsf{RT}_\mathsf{der} = \emptyset$. An item is not divisible and cannot be merged or converted.
	\item \textbf{Issuance policies}: open issuance for a network-wide type in which anyone mints; authorized issuance for a curated collection, in which the type is the collection and its issuer authorizes each item; capped profiles V (default), S (count mode) or P for a collection of bounded size; proof-gated issuance for type-defined rules.
	\item \textbf{Metadata}: for a collection type, a registry entry with $\texttt{assetKind} = \texttt{"non-fungible"}$ describes the collection (Sec.~\ref{app:asset-metadata}); per-item metadata is the payload.
\end{itemize}

Transfer-level auxiliary data is free-form transport (Sec.~\ref{sec:mint-authorization}) and never redefines the item.

\subsubsection{Item Structures}\label{app:nft-items}

\index{NFT item}
An \emph{NFT item} is one of four structures, each a CBOR array under its own semantic tag. The first element is the format version, $1$; the arity is fixed per version, so adding a field requires a new version.

\begin{longtable}{llp{9.2cm}}
\textbf{Tag} & \textbf{Name} & \textbf{Structure} \\
\hline
\endhead
39052 & NftMetadata & $(1, \mathsf{name}, \mathsf{description}, \mathsf{image}, \mathsf{animation},$\newline\hspace*{1em}$\mathsf{url}, \mathsf{attributes}, \mathsf{collection}, \mathsf{cid})$: item description \\
39053 & NftMedia & $(1, \mathsf{mt}, b)$: a file $b$ of media type $\mathsf{mt}$, carried inline \\
39054 & NftLink & $(1, \mathsf{mt}, \mathsf{uri}, h)$: a file of media type $\mathsf{mt}$, hosted at $\mathsf{uri}$, with SHA-256 digest $h$ \\
39055 & NftSigned & $(1, \mathsf{pk}_\mathsf{c}, c, \sigma)$: an item $c$ with the creator's key and signature \\
\end{longtable}

NftMetadata fields:

\begin{longtable}{c>{\raggedright\arraybackslash}p{3.1cm}>{\raggedright\arraybackslash}p{5.6cm}l}
\textbf{Index} & \textbf{Type} & \textbf{Description} & \textbf{ERC-721} \\
\hline
\endhead
0 & unsigned integer & format version, $1$ & \\
1 & text & item name; required & \texttt{name} \\
2 & text or \texttt{null} & description & \texttt{description} \\
3 & media reference or \texttt{null} & image & \texttt{image} \\
4 & media reference or \texttt{null} & video, audio or 3D content & \texttt{animation\_url} \\
5 & text or \texttt{null} & external web page of the item & \texttt{external\_url} \\
6 & array & traits: pairs $(\mathsf{trait\_type}, \mathsf{value})$, with text $\mathsf{trait\_type}$ and text or integer $\mathsf{value}$; empty when there are none & \texttt{attributes} \\
7 & text or \texttt{null} & collection display name (Sec.~\ref{app:nft-collection}) & contract \texttt{name} \\
8 & byte string or \texttt{null} & claimed collection identifier $\mathsf{cid}$, $1 \le \len{\mathsf{cid}} \le 64$ (Sec.~\ref{app:nft-collection}) & contract address \\
\end{longtable}

\noindent A \emph{media reference} is an NftMedia item or an NftLink item that is not a document link (below). The remaining structures carry:
\begin{description}
	\item[NftMedia] $\mathsf{mt}$, a media type in the sense of RFC~6838 (\texttt{type/subtype}, lower case, without parameters); $b$, the non-empty file.
	\item[NftLink] $\mathsf{mt}$, the media type of the linked file; $\mathsf{uri}$, a locator with scheme \texttt{https}, \texttt{ipfs} or \texttt{ar}; $h$, the 32-byte SHA-256 digest of the file's bytes. The digest authenticates the content and the locator is a non-authoritative hint. An IPFS content identifier is not a substitute for $h$, since for most files it is not a hash of the raw bytes.
	\item[NftSigned] $\mathsf{pk}_\mathsf{c}$, a 33-byte compressed secp256k1 public key; $c$, an embedded item (a CBOR item, not a byte string); $\sigma$, a 65-byte signature $R \mathbin{\|} S \mathbin{\|} V$ in the format of Appendix~\ref{app:signature-format}.
\end{description}

Items are encoded in deterministic CBOR (Appendix~\ref{app:encodings}) with definite lengths. Maps, floating-point numbers, booleans and simple values other than \texttt{null} do not occur. Text is valid UTF-8. An optional text field is either \texttt{null} or non-empty, so absence has a single encoding. Integer trait values $v$ satisfy $\len{v} \le 2^{53}-1$; fractional or larger numbers are written as text. The integer range is the range that common client languages represent exactly. The ERC-721 fields \texttt{background\_color} and \texttt{youtube\_url} and the trait fields \texttt{display\_type} and \texttt{max\_value} are omitted.

\paragraph{Document links.} An NftLink with $\mathsf{mt} = \texttt{application/vnd.unicity.nft+cbor}$ is a \emph{document link}. It references an NftMetadata or NftMedia item hosted outside the token, in place of a media file; it is the analogue of an ERC-721 \texttt{tokenURI} that points at a metadata document. The linked file is exactly one encoded item. It is never a link, so links do not chain, and it is never an NftSigned item, since a signature binds one token and therefore stays in the token.

\paragraph{Placement.} Items nest as follows; an item in any other position makes the enclosing payload or document unrecognized.

\begin{longtable}{>{\raggedright\arraybackslash}p{5.4cm}>{\raggedright\arraybackslash}p{8.8cm}}
\textbf{Position} & \textbf{Admissible items} \\
\hline
\endhead
genesis payload $\auxd'_0$ & NftMetadata, NftMedia, NftLink (including a document link), NftSigned \\
the item $c$ of an NftSigned & NftMetadata, NftMedia, NftLink (including a document link) \\
$\mathsf{image}$, $\mathsf{animation}$ & media reference or \texttt{null} \\
file referenced by a document link & NftMetadata, NftMedia \\
\end{longtable}

\subsubsection{Recognition and Validity}\label{app:nft-validity}

Let $\mathcal{N}$ be the set of well-formed items admissible at the top level. The \emph{recognition} function $\Call{RecognizeNFT}{b} \in \mathcal{N} \opt$ returns $c$ if the byte string $b$ is the complete encoding of one item $c \in \mathcal{N}$ satisfying the rules of Sec.~\ref{app:nft-items}, and $\bot$ otherwise. An absent payload, raw file bytes, an untagged array, an unknown tag or version, a wrong arity or field type, a non-deterministic encoding and trailing bytes all yield $\bot$. For a UTS-2 type, $\Call{ParseAux}{\mathsf{ty}, \auxd'}$ returns $c = \Call{RecognizeNFT}{\auxd'}$ if $c \ne \bot$ and $c$ satisfies the collection rule of Sec.~\ref{app:nft-collection}, and $\bot$ otherwise.

In a UTS-2 type the item format is therefore a validity condition, checked by Alg.~\ref{alg:mint-auth} like any other type rule: a token whose payload is not a recognized item is rejected. This differentiates UTS-2 from an open container, whose wallets display recognized items (Sec.~\ref{app:uts-blank}) without specific validation.

Validity depends only on verifier-held data. none of the following will invalidate a UTS-2 token:
\begin{itemize}
	\item \emph{Linked content.} A reader fetches a linked file, requires its SHA-256 digest to equal $h$ before rendering it, and for a document link additionally requires the file to be a recognized item. A file that is unavailable, fails the digest check or does not parse is not displayed. The token remains valid and transferable.
	\item \emph{Creator attribution.} The attribution status of Sec.~\ref{app:nft-attribution} is computable from the token, but it is a statement about authorship, not about authorization to issue. A type that requires a particular signer expresses that requirement as its issuance policy.
	\item \emph{Presentation policy.} Wallets render text as plain text, render only media types on an allowlist, and may limit the size of inline media. These are presentation rules.
\end{itemize}

The genesis payload of a token does not change, thus a recognition result, including $\bot$, may be cached for its token identifier indefinitely. Inline media is part of token and stored, transferred together; hosted media helps to reduce token size.

\subsubsection{Creator Attribution}\label{app:nft-attribution}

\index{NFT item!creator attribution}
An NftSigned item binds its embedded item $c$ to one token and one first owner. Define $\mathsf{NFT\_SIGNED}$ as the ASCII byte string \texttt{"UNICITY\_NFT\_SIGNED"}. For mint data $D_\mathsf{mint} = (\alpha, \predi', \mathsf{salt}, \mathsf{ty}, \mathsf{reason}, \auxd')$ whose payload is an NftSigned item $(1, \mathsf{pk}_\mathsf{c}, c, \sigma)$, let $b_c$ be the bytes of $c$ exactly as they occur in $\auxd'$, tag included. The \emph{attribution digest} is
\[
m_\mathsf{nft} := \mathsf{SHA\text{-}256}\bigl(\mathsf{CBOR}(
	\mathsf{NFT\_SIGNED},\;
	1,\;
	\alpha,\;
	\encpred(\predi'),\;
	\mathsf{id},\;
	\mathsf{ty},\;
	b_c
)\bigr),
\]
with $\mathsf{id} = H(\mathsf{salt}, \alpha)$. The format version and $\alpha$ are unsigned integers and every other term is a byte string, following the layout of the genesis commitment (Sec.~\ref{app:genesis-commitment}). The signature $\sigma$ lies outside $b_c$, so the digest does not cover itself. A verifier hashes the bytes it received and does not re-encode the decoded item. The \emph{attribution status} $\mathsf{attr}(T_0) \in \{\mathsf{unsigned}, \mathsf{valid}, \mathsf{invalid}\}$ of a genesis is $\mathsf{unsigned}$ if its payload is not an NftSigned item. Otherwise it is $\mathsf{valid}$ if all of the following hold, and $\mathsf{invalid}$ if any fails:
\begin{enumerate}
	\item $\mathsf{pk}_\mathsf{c}$ is a valid compressed secp256k1 point;
	\item $\sigma = R \mathbin{\|} S \mathbin{\|} V$ with $1 \le S \le n/2$, where $n$ is the order of the secp256k1 group, and $V \in \{0, \ldots, 3\}$;
	\item public-key recovery from $(R, S, V)$ over $m_\mathsf{nft}$ yields $\mathsf{pk}_\mathsf{c}$.
\end{enumerate}
The recovery identifier is bound, so the same $(R, S)$ under a different $V$ is invalid, and the low-$S$ requirement excludes the second, malleated form of a valid signature. The digest uses the genesis recipient $\predi'$, not the current owner, so the status of a token does not change under transfer.

The terms of $m_\mathsf{nft}$ exclude the following misuses:
\begin{description}
	\item[Domain separator and version] reuse of the signature in any other protocol. $\mathsf{pk}_\mathsf{c}$ is typically a wallet key that also signs transactions and messages.
	\item[Token identifier] copying a signed item onto another token.
	\item[Genesis recipient] theft by front-running. An observer of the mint request can resubmit it with the same salt, hence the same identifier, and itself as recipient; the resulting token's status is $\mathsf{invalid}$. The observer can still consume the identifier first and block the creator's mint. That is denial of issuance, not theft: the creator mints again with a fresh salt.
	\item[Token type] minting the signed item under another type. The identifier $H(\mathsf{salt}, \alpha)$ does not commit to $\mathsf{ty}$, so the same salt yields the same identifier under every type.
	\item[Network identifier] leaving the network implicit. The identifier already commits to $\alpha$, so this term states the network explicitly rather than adding a replay defense.
\end{description}

\paragraph{Attribution is not authorization.} A $\mathsf{valid}$ status establishes that the holder of the secret key of $\mathsf{pk}_\mathsf{c}$ signed $c$ for this token and this first owner. It establishes neither that $\mathsf{pk}_\mathsf{c}$ belongs to a recognized creator, which is a presentation-level lookup (e.g.\ to a name tag), nor that issuance was authorized. Authorization is the mint reason, verified under the type's issuance policy. The two statements differ even when one party makes both, and they are signed over different digests with different domain separators ($m_\mathsf{mint}$ and $m_\mathsf{nft}$). Under open issuance anyone can mint an item that names a well-known key together with an arbitrary signature. A wallet therefore MUST NOT attribute a token to $\mathsf{pk}_\mathsf{c}$ unless the status is $\mathsf{valid}$, and SHOULD present a valid status as \emph{signed}, not as verified. It SHOULD keep three claims apart: that the signature is valid, that the signer is a key the wallet recognizes, and that the item belongs to a collection.

An NftSigned item over a document link signs the link including $h$, and therefore the document. A payload that is pre-committed before its mint (profile S, Sec.~\ref{app:uts-slots}) can carry an attribution only if the first owner is fixed at commitment time, since $m_\mathsf{nft}$ binds $\predi'$.

\subsubsection{Collection Identity}\label{app:nft-collection}

\index{NFT collection}
ERC-721 identifies a collection by the address of its contract, and the contract holds the authority to mint. Here the analogue of the contract is the token type. Under authorized, capped or proof-gated issuance, $\mathsf{ty}$ is the collection: its manifest fixes the issuance rules, and a valid token proves membership by its verified reason. Under open issuance the type is shared by unrelated minters and identifies no collection.

The payload fields $\mathsf{collection}$ and $\mathsf{cid}$ are claims of the minter:
\begin{itemize}
	\item $\mathsf{collection}$ is a display name. Names collide and never identify a collection.
	\item For a type under any policy other than open issuance, $\Call{ParseAux}{}$ requires $\mathsf{cid} \in \{\bot, \mathsf{ty}\}$, so the claimed identity cannot contradict the authenticated one. A payload committed in the manifest itself (profile S) cannot contain $\mathsf{ty}$ and uses $\bot$. A wallet takes the collection of such a token from $\mathsf{ty}$.
	\item Under open issuance, $\mathsf{cid}$ is unrestricted, e.g.\ the type identifier of a collection or the hash of a collection manifest. It is not verified, and a wallet MUST NOT present it as verified membership, whether or not the item is signed.
\end{itemize}

Collection-level metadata (name, description, icon) is the registry entry of the collection type, the analogue of contract-level metadata. The item format defines no supply limit, and no field limits the number of items. A bound on the size of a collection is a property of the type's issuance policy (Sec.~\ref{app:uts-capped}). Tokens with distinct identifiers may carry byte-identical items. Under authorized issuance the issuer is responsible for not authorizing duplicates; under open issuance nothing prevents them (Sec.~\ref{app:policy-open}).

\paragraph{Serial numbers.} An ERC-721 contract chooses each \texttt{tokenId}, often sequentially. Here the item's identity is $\mathsf{id}$, derived from the minter's salt. The capped profiles V, S and P additionally fix a serial $j$ in the verified reason, and that serial, not any payload field, is the item's position in the collection. Open and authorized issuance define no serials; a name such as ``Cool Cat \#1'' is display data.

\begin{longtable}{>{\raggedright\arraybackslash}p{4.6cm}>{\raggedright\arraybackslash}p{9.6cm}}
\textbf{ERC-721} & \textbf{UTS-2} \\
\hline
\endhead
contract address & token type $\mathsf{ty}$ of a collection type; none under open issuance \\
\texttt{tokenId} & token identifier $\mathsf{id}$; the serial $j$ under a capped profile \\
\texttt{ownerOf} & current owner predicate, known to the holder; no public lookup \\
\texttt{balanceOf} & number of items a wallet holds, computed by the holder \\
\texttt{transferFrom}, \texttt{safeTransferFrom} & transfer transaction; the receiver-side check of \texttt{safeTransferFrom} corresponds to the recipient's verification of the token before acceptance \\
\texttt{approve}, \texttt{setApprovalForAll} & no approval registry; delegation is expressed by the owner predicate (Sec.~\ref{sec:standard-predicates}) \\
\texttt{name}, \texttt{symbol} & registry metadata of the collection type \\
\texttt{tokenURI} and metadata document & genesis payload item; a document link for hosted metadata \\
\texttt{totalSupply}, \texttt{tokenByIndex} (enumeration extension) & no enumeration; capped issuance bounds the number of items, not the number issued (Sec.~\ref{app:capped-availability}) \\
mint and burn & direct issuance under the type's policy; burn transaction (Sec.~\ref{sec:burn-transaction}) \\
\end{longtable}

\subsection{UTS-3: Pooled Assets}\label{app:uts-pool}

\index{token pool}
A \emph{pool} aggregates value of one fungible asset type into a single token under shared control. With the merge reason (Sec.~\ref{app:reason-merge}) as the aggregation primitive, two profiles are distinguished by where the pool's accounting logic runs.

\paragraph{Profile A: Committee pool.} Fully expressible with the machinery already specified:
\begin{itemize}
	\item The \emph{pool token} is an ordinary UTS-1 token whose owner predicate is a multi-party condition: $\predi_\mathsf{msig}$, $\predi_\mathsf{tsig}$ ($k$-of-$n$), or a P2SH-wrapped governance predicate (Sec.~\ref{sec:standard-predicates}). The underlying asset type includes $\mathsf{MERGE\_TAG}$ and 39044 in $\mathsf{RT}$.
	\item \emph{Deposit}: the depositor transfers a token of the pooled type to the pool predicate; the committee merges it into the pool token. \emph{Withdrawal}: the committee splits the pool token, sending one output to the withdrawer and the remainder back to the pool predicate.
	\item \emph{Shares}, where needed, are a separate UTS-1 type under authorized issuance whose issuance predicate is the same committee: deposits are acknowledged by share issuance, withdrawals demand share burn.
\end{itemize}
Custody of the pooled value is protected cryptographically by the threshold predicate and by value conservation of merge and split --- the committee cannot inflate the pool. What rests on the committee's honesty is \emph{share accounting}: issuing shares consistently with deposits and honoring redemptions. This profile therefore suits treasuries, escrow and OTC pools where the committee is the counterparty anyway.

\paragraph{Profile B: Autonomous pool.} Removing the committee means enforcing the pool's accounting invariant by verification instead of trust: the pool token becomes a type-governed state machine, with bearer share tokens minted and redeemed against its certified transitions. This profile is specified in Sec.~\ref{app:autonomous-pools}.

\section{Autonomous Pools}\label{app:autonomous-pools}

\index{token pool!autonomous}
An \emph{autonomous pool} aggregates value of one fungible asset under no party's control: the pool's accounting invariant is enforced by type-defined validation rules that every verifier re-runs, not by a committee. This section resolves the three problems left open by the committee profile (Sec.~\ref{app:uts-pool}): cross-type value flow is handled by the conversion format (Sec.~\ref{app:reason-convert}); the share ledger is externalized into bearer share tokens, keeping the pool state constant-size; and deposit replay is excluded by a nullifier accumulator carried in the pool state, borrowing the mechanism of Sec.~\ref{sec:bridge-replay}.

\subsection{Setting}

Three token types cooperate; all three must be registered in a verifier's $\Gamma.\mathcal{TD}$ for pool-derived tokens to verify (fail-shut otherwise):
\begin{description}
	\item[Underlying asset $\mathsf{ty}_\mathsf{A}$] a fungible type with UTS-1 payload rules whose manifest admits $\mathsf{CONVERT\_TAG}$ with $\mathsf{ty}_\mathsf{P} \in \mathcal{CS}$. The asset type opts in to the pool by its frozen definition, since its own verifiers must be able to re-check the pool's conservation. For a capped asset, this requires a separately specified extension of capped issuance (Sec.~\ref{app:uts-capped}): the standard profiles exclude conversion, and the extension must preserve the cap across plain and pooled value using the conservation argument below.
	\item[Pool type $\mathsf{ty}_\mathsf{P}$] the state machine defined below. One \emph{pool instance} is one token of this type; instances are distinguished by their token identifiers $\mathsf{id}_\mathsf{P}$, and anyone may open one ($\mathsf{RT} = \{\bot\}$).
	\item[Share type $\mathsf{ty}_\mathsf{S}$] a fungible type with $\mathsf{RT} = \{\mathsf{CONVERT\_TAG}, 39044, \mathsf{MERGE\_TAG}\}$ and $\mathcal{CS} = \{\mathsf{ty}_\mathsf{P}\}$. Shares use the asset framework (Sec.~\ref{sec:asset-structure}) with \emph{asset identifier $\mathsf{aid} = \mathsf{id}_\mathsf{P}$}: shares of different pool instances are distinct assets, so split and merge conservation automatically prevents mixing them. Shares are ordinary bearer tokens --- transferable, splittable, mergeable.
\end{description}

The pool token is \emph{protocol-owned}: its genesis predicate and, by transfer rule, every subsequent owner predicate is the open predicate $\predi_\mathsf{open}$ (Sec.~\ref{sec:open-predicate}). Anyone can spend the pool state, but only along transitions valid under $\mathsf{ty}_\mathsf{P}$'s rules, and never away from $\predi_\mathsf{open}$ --- the pool cannot be taken private.

\subsection{Pool State}

The pool state is the tuple
\[
\sigma = (\mathsf{aid},\; V,\; S,\; A),
\]
where $\mathsf{aid}$ is the pooled asset identifier, $V$ the pooled value, $S$ the shares outstanding, and $A$ the root of the \emph{consumed-burn accumulator}: a sparse Merkle tree (Sec.~\ref{sec:smt}) over the state identifiers of every deposit and withdrawal burn the pool has absorbed. The genesis payload is $\sigma_0 = (\mathsf{aid}, 0, 0, A_\varnothing)$ with $\mathsf{reason} = \bot$; each transfer's auxiliary data carries a \emph{transition record} and the successor state $\sigma'$; the current pool state is the last certified $\sigma$. All amounts satisfy $0 \le V, S < 2^{256}$, with overflow rejected.

\subsection{Transitions}

A transition is a transfer of the pool token. Checks common to every transition (run as $\mathsf{ty}_\mathsf{P}$'s per-transfer check, Sec.~\ref{sec:mint-authorization}, step~6 of Sec.~\ref{sec:transfer-tx}):
\begin{enumerate}
	\item the successor predicate is exactly $\predi_\mathsf{open}$;
	\item the auxiliary data decodes strictly to $(\mathsf{record}, \sigma')$, and $\sigma'$ equals the state computed from the predecessor $\sigma$ and the record per the rules below;
	\item the record's nullifier is proven absent from $A$ and $\sigma'.A$ is the root after inserting it (the SMT provides both certificates, cf.\ Sec.~\ref{sec:bridge-replay}).
\end{enumerate}

\paragraph{Deposit.} The depositor burns a token of type $\mathsf{ty}_\mathsf{A}$ carrying $(\mathsf{aid}, v)$, with burn manifest (under tag $\mathsf{POOL\_DEPOSIT\_TAG}$) committing $(\mathsf{id}_\mathsf{P},\; h_\mathsf{claim})$, where $h_\mathsf{claim} = (H(\mathsf{CBOR}(\predi'_\mathsf{s})), \mathsf{salt}_\mathsf{s})$ fixes the future share token's recipient and salt. The deposit record embeds the burned token (verified via the collector, Sec.~\ref{sec:nested-verification}) and the transition check requires:
\begin{enumerate}
	\item the embedded token is of type $\mathsf{ty}_\mathsf{A}$, network $\mathcal{T}.\alpha$, its final transfer is a burn binding the manifest above with this pool's $\mathsf{id}_\mathsf{P}$, and its value for $\mathsf{aid}$ is $v \ge 1$;
	\item the share amount is
	\[
	\Delta s =
	\begin{cases}
		v & \text{if } S = 0\\
		\lfloor v \cdot S / V \rfloor & \text{otherwise,}
	\end{cases}
	\qquad \Delta s \ge 1 \text{ (dust deposits rejected)};
	\]
	\item $V' = V + v$, $S' = S + \Delta s$;
	\item the record contains the conversion output commitment $d_\mathsf{s}$ (Sec.~\ref{app:reason-convert}) for a $\mathsf{ty}_\mathsf{S}$ mint with payload $\{(\mathsf{id}_\mathsf{P}, \Delta s)\}$, recipient predicate and salt as fixed by $h_\mathsf{claim}$;
	\item the nullifier is the burned token's final state identifier.
\end{enumerate}
Since the burn commits $h_\mathsf{claim}$, a third party who includes someone else's burn in a transition can only execute that deposit as intended --- the shares are creditable solely to the committed recipient. Share amounts depend on $(V, S)$ at execution time, which is why $\Delta s$ and $d_\mathsf{s}$ are computed in the transition record, not in the burn.

\paragraph{Withdrawal.} The holder burns share tokens carrying $(\mathsf{id}_\mathsf{P}, s)$, $1 \le s \le S$, with burn manifest (under tag $\mathsf{POOL\_WITHDRAW\_TAG}$) committing $(\mathsf{id}_\mathsf{P},\; h_\mathsf{rcpt})$ for the withdrawn token's recipient and salt. The withdrawal record embeds the share burn, and the transition check requires:
\begin{enumerate}
	\item the embedded token is of type $\mathsf{ty}_\mathsf{S}$ with asset identifier exactly $\mathsf{id}_\mathsf{P}$, burned as above;
	\item the withdrawn value is $w = \lfloor s \cdot V / S \rfloor$ with $w \ge 1$;
	\item $V' = V - w$, $S' = S - s$;
	\item the record contains the conversion output commitment $d_\mathsf{w}$ for a $\mathsf{ty}_\mathsf{A}$ mint with payload $\{(\mathsf{aid}, w)\}$, recipient and salt per $h_\mathsf{rcpt}$;
	\item the nullifier is the share burn's final state identifier.
\end{enumerate}
Rounding is always downward: dust remains in the pool, to the benefit of remaining share holders. The floor arithmetic preserves the invariant $S > 0 \Rightarrow V > 0$: a partial withdrawal ($s < S$) leaves $w \le \lfloor (S-1)V/S \rfloor < V$, and only a full withdrawal ($s = S$) can zero the pool.

\subsection{Claiming Outputs}

Shares and withdrawn tokens are minted with the conversion reason (Sec.~\ref{app:reason-convert}): the reason embeds the pool token up to and including the committing transition, and the verifier matches the recomputed output commitment against the record. Each commitment authorizes exactly one mint (the committed salt fixes the output's genesis state, certified at most once). An output left unclaimed is its claimant's loss alone --- the pool state already accounts for it, and no other party's value is affected, consistent with the platform's uniform failure direction: value can be destroyed by negligence, never created.

\subsection{Security}

\begin{description}
	\item[Solvency] Every unit of $V$ entered through a certified $\mathsf{ty}_\mathsf{A}$ burn counted exactly once (the accumulator rejects replayed burns); every exit is a $V$-decrement with $w \le V$ by construction. The pool can never commit to releasing more than was deposited.
	\item[Cross-type conservation] The total supply of $\mathsf{ty}_\mathsf{A}$ --- plain tokens plus $\sum_\mathsf{pools} V$ --- is invariant under pool operations: deposits burn plain value (single-spend), withdrawals mint plain value only against a $V$-decrement, and $V$ increments only against burns. Every verifier of a withdrawal-minted token re-verifies the entire pool history through the worklist, so this is not an assumption but a checked property.
	\item[Pro-rata fairness] Share value $V/S$ is non-decreasing under deposits and withdrawals (floor rounding); no transition dilutes existing holders in favor of a newcomer.
	\item[No capture] Transitions force $\predi' = \predi_\mathsf{open}$ and the state equations; there is no code path by which any party, including the pool's creator, spends pool value except via a withdrawal backed by share burn.
	\item[Race safety] Two parties spending the same pool state race; single-spend lets exactly one transition certify. The loser's burn remains valid --- it binds $\mathsf{id}_\mathsf{P}$, not a specific pool state --- and is simply resubmitted in a subsequent transition.
\end{description}

\subsection{Two-Asset Pools and Swaps}\label{app:pool-swaps}

\index{token pool!swap}\index{automated market maker}
A \emph{two-asset pool} extends the machinery above to an automated market maker: liquidity providers deposit two assets, swappers exchange one for the other against the pooled reserves, and the price is set by the reserve ratio. The state generalizes to
\[
\sigma = (\mathsf{aid}_1,\; V_1,\; \mathsf{aid}_2,\; V_2,\; S,\; A),
\qquad \mathsf{aid}_1 < \mathsf{aid}_2 \text{ (canonical order)},
\]
with both assets carried by the same underlying type $\mathsf{ty}_\mathsf{A}$ (a fungible payload holds multiple assets, Sec.~\ref{sec:asset-structure}). The manifest of $\mathsf{ty}_\mathsf{P}$ additionally pins a \emph{fee} $\varphi$ in per-mille, $0 \le \varphi < 1000$. The common transition checks (open successor predicate, state derivation, nullifier accumulator) are unchanged; the transitions are:

\paragraph{Liquidity deposit.} The provider burns a $\mathsf{ty}_\mathsf{A}$ token carrying $\{(\mathsf{aid}_1, v_1), (\mathsf{aid}_2, v_2)\}$, $v_1, v_2 \ge 1$, with a deposit manifest as before, optionally extended with a minimum share amount $\Delta s_\mathsf{min}$. The transition requires
\[
\Delta s =
\begin{cases}
	\lfloor\sqrt{v_1 \cdot v_2}\rfloor & \text{if } S = 0\\
	\min\!\left(\lfloor v_1 S / V_1 \rfloor,\; \lfloor v_2 S / V_2 \rfloor\right) & \text{otherwise,}
\end{cases}
\qquad \Delta s \ge \max(1, \Delta s_\mathsf{min}),
\]
and $V_i' = V_i + v_i$, $S' = S + \Delta s$. The $\min$ rule donates any contribution in excess of the current reserve ratio to the pool, so providers match the ratio in their own interest; $\Delta s_\mathsf{min}$ bounds the loss if reserves move between burn and execution.

\paragraph{Liquidity withdrawal.} As in the single-asset pool, per asset: $w_i = \lfloor s \cdot V_i / S \rfloor$, $V_i' = V_i - w_i$, $S' = S - s$; the committed output payload is $\{(\mathsf{aid}_1, w_1), (\mathsf{aid}_2, w_2)\}$ with zero-valued entries omitted and at least one entry positive.

\paragraph{Swap.} The swapper burns a $\mathsf{ty}_\mathsf{A}$ token carrying $\{(\mathsf{aid}_\mathsf{in}, v_\mathsf{in})\}$, with a swap manifest (tag $\mathsf{POOL\_SWAP\_TAG}$) committing $(\mathsf{id}_\mathsf{P}, \mathsf{aid}_\mathsf{in}, h_\mathsf{rcpt}, v_\mathsf{min}, h_\mathsf{refund})$: the pool instance, the input asset, the output's recipient commitment, a minimum acceptable output, and a refund recipient commitment. The transition requires $S > 0$,
\[
v_\mathsf{out} \;=\; \left\lfloor
\frac{v_\mathsf{in} \, (1000 - \varphi) \, V_\mathsf{out}}
     {1000 \, V_\mathsf{in} + (1000 - \varphi) \, v_\mathsf{in}}
\right\rfloor,
\qquad
v_\mathsf{out} \ge \max(1, v_\mathsf{min}),
\]
$V_\mathsf{in}' = V_\mathsf{in} + v_\mathsf{in}$, $V_\mathsf{out}' = V_\mathsf{out} - v_\mathsf{out}$, and the record commits a $\mathsf{ty}_\mathsf{A}$ output $\{(\mathsf{aid}_\mathsf{out}, v_\mathsf{out})\}$ per $h_\mathsf{rcpt}$. The formula is the constant-product rule with fee: it implies $V_1' \cdot V_2' \ge V_1 \cdot V_2$, with strict growth for $\varphi > 0$ --- the reserve product never decreases, so swaps can only add value per share.

\paragraph{Refund.} Any pool-bound burn of this profile (deposit or swap) may instead be consumed by a \emph{refund} transition: the pool state is unchanged except for the nullifier insertion, and the record commits a $\mathsf{ty}_\mathsf{A}$ output returning the burn's exact asset content to the manifest's $h_\mathsf{refund}$. This guarantees that a burn whose bound ($\Delta s_\mathsf{min}$, $v_\mathsf{min}$) has become unsatisfiable is never stranded: its value is always recoverable. Whoever executes the transition chooses between execution and refund; a hostile executor can therefore deny a trade but never confiscate --- the refund pays only the committed refund recipient.

\paragraph{Security notes.} Solvency, cross-type conservation, no-capture and race safety carry over unchanged; the accumulator prevents replaying any burn across transitions. Two properties are specific to swaps:
\begin{description}
	\item[Provider protection] The reserve-product invariant is checked by every verifier per transition; no sequence of swaps, at any prices, dilutes the value backing a share below the constant-product curve.
	\item[Ordering exposure] Whoever wins the race for the pool state orders the pending trades; sandwich-style reordering is possible as in any automated market maker. The swapper's $v_\mathsf{min}$ caps the damage per trade, and the refund path bounds the downside to a cancelled trade.
\end{description}

\subsection{Limitations and Open Items}

\begin{itemize}
	\item \emph{Throughput}: one spendable state serializes all operations on a pool instance; contention is inherent. Sharded or tree-structured pools are a candidate extension.
	\item \emph{Verification cost and compression}: pool verification cost grows with transition count, and each deposit/withdrawal embeds a burned token. Practical pools depend on history compression of $\mathsf{ty}_\mathsf{P}$ tokens. Since the pool type defines per-transfer checks, the generic chain relation $\mathfrak{R}_\mathsf{ch}$ (Sec.~\ref{sec:compression-relations}) does not cover it: a variant carrying the pool transition checks is required, and $\mathsf{ty}_\mathsf{P}$'s manifest must pin its verification keys --- tracked in the roadmap.
	\item \emph{Privacy}: the pool state and every transition are visible to all verifiers of pool-derived tokens; deposits, swaps and withdrawals are linkable to their claim commitments.
	\item \emph{Encodings}: transition record and manifest encodings, and the accumulator's tree parameters, are fixed when this draft freezes.
\end{itemize}

\section{Specification Roadmap}\label{app:staging}

This appendix is written iteratively, with review feedback between revisions. Each item carries a document status: \emph{specified} --- normative and stable; \emph{draft} --- complete but subject to change until reviewed and frozen; \emph{outline} --- design direction with open questions, to be elaborated in a later revision. The roadmap concerns this document only; implementation and deployment planning are tracked separately.

\begin{longtable}{p{7.2cm}lp{4.6cm}}
\textbf{Item} & \textbf{Status} & \textbf{Next revision step} \\
\hline
\endhead
Mint authorization flow, verifier registry, fail-shut policy (Sec.~\ref{sec:token-types}) & specified & --- \\
Split reason, bridge backing reason (Appendices~\ref{app:split-proof-encoding}, \ref{app:bridging}) & specified & --- \\
Registries and record schemas (Sec.~\ref{app:registries}--\ref{app:type-registry}) & draft & freeze record schema after review \\
Genesis commitment, issuer-signature reason (Sec.~\ref{app:genesis-commitment}, \ref{app:reason-issuer}) & draft & review; allocate tags \\
Committed-slot reason: count and value modes (Sec.~\ref{app:reason-slot}) & draft & review; allocate tag \\
VRF and public serial reasons (Secs.~\ref{app:reason-vrf}, \ref{app:reason-serial}) & draft & review; allocate tags; interoperability vectors for serial derivation \\
Merge reason (Sec.~\ref{app:reason-merge}) & draft & review; allocate tags \\
Conversion reason (Sec.~\ref{app:reason-convert}) & draft & review; allocate tags \\
Proof-carrying reason (Sec.~\ref{app:reason-proof}) & draft & select proof-system profile \\
History compression (Sec.~\ref{sec:history-compression}) & draft & select proof-system profile; encode $\mathcal{H}$; relation variant for per-transfer type checks \\
Provenance compression: token evidence, certificates, $\mathfrak{R}_\mathsf{lin}$ (Sec.~\ref{sec:provenance-compression}) & draft & review; encode $\mathcal{K}$ and the evidence set; reference form for conversion and pool records \\
Reference-form split reason (Sec.~\ref{app:reason-splitref}) & draft & review; allocate tag \\
Issuance policies: open, authorized, bridged, proof-gated (Sec.~\ref{app:issuance-policies}) & draft & review \\
UTS-0, UTS-1 (Secs.~\ref{app:uts-blank}, \ref{app:uts-fungible}) & draft & review \\
UTS-2 item format, recognition, creator attribution, collection identity (Sec.~\ref{app:uts-nft}) & draft & review; reconcile $\encpred$ with the tagged predicate encoding (tag 39032) used by the implementation, which affects the genesis, split and attribution commitments \\
Capped issuance: singleton, VRF, committed-slot and public-serial profiles (Sec.~\ref{app:uts-capped}) & draft & review profile policies and cap/availability separation \\
UTS-3 pooled assets: committee profile (Sec.~\ref{app:uts-pool}) & draft & review \\
Autonomous pools, incl.\ two-asset swaps (Sec.~\ref{app:autonomous-pools}) & draft & review; freeze record encodings \\
\end{longtable}

Later revisions extend this appendix; a revision MUST NOT alter the verification semantics of any item already frozen as specified.
